% !TEX TS-program = xelatex
% !TEX encoding = UTF-8
%======================================================================
%  SDR Lab -- Experiment 2: FM radio, demodulation, stereo and RDS
%  LAB MANUAL for BOTH sessions (weeks 3 and 4 of the regulations)
%  File : labs/lab_02_fm.tex                 (XeLaTeX, compile twice)
%
%  [INSTRUCTOR] QUICK FACTS (details: instructor/instructor_notes_02_fm.md)
%   * Session 1 = Parts 1-5 (two detectors, cabled NBFM link, threshold,
%     team challenge "own deviation").  Session 2 = Parts 6-9 (WFM mono,
%     stereo, RDS, challenge "RDS versus level").  The second challenge is
%     my design: the regulations list one team challenge for this experiment
%     but the session pattern has a challenge slot in every session.
%   * Receiver: RTL-SDR.  Transmitter: written for ADALM-PLUTO ("PlutoSDR
%     Sink"); Table tab:txhw maps the parameters to HackRF / USRP.  [HW]
%   * Transmission is CABLED ONLY: TX port -> 30 dB attenuator -> receiver.
%     You must hand out the list of allowed carrier frequencies (f_tx).
%     A stock Pluto tunes 325 MHz - 3.8 GHz, so the default f_tx is 434 MHz.
%   * All six flowgraphs exist as reference files in instructor/lab02/ and
%     compile with grcc (GNU Radio 3.10.9).  Both simulations were run; the
%     hardware flowgraphs were NOT run (no hardware in my test environment).
%   * Session 2 needs gr-rds on the lab PCs and the two helper files
%     code/lab02/sdrlab_rds_counter.py and instructor/lab02/lab02_stereo_rds_tx.grc
%     (the latter is handed to the students as a ready-made signal source).
%   * Numbers quoted in "expected observation" boxes come from my GNU Radio
%     simulations with ideal noise; real hardware will differ by a few dB.
%======================================================================
\documentclass[a4paper,11pt]{article}
\IfFileExists{../common/preamble.tex}{\def\SDRroot{..}}{\def\SDRroot{.}}
\def\SDRdoctype{lab}
\def\SDRexpnum{2}
\def\SDRexptitle{رادیو \lr{FM}}
% [orig] \def\SDRexptitle{رادیو \lr{FM}، استریو و \lr{RDS}}
\def\SDRexptitleEN{FM Radio}
% [orig] \def\SDRexptitleEN{FM Radio, Stereo and RDS}
\input{\SDRroot/common/preamble}

% banner artwork: stereo multiplex spectrum (drawn, not a screenshot)
\renewcommand{\SDRbannerart}{%
  \draw[sdrLight!40,line width=0.4pt] (0.2,0.45)--(4.0,0.45);
  \fill[sdrBlue!70] (0.3,0.45)--(0.4,1.75)--(1.1,1.75)--(1.2,0.45)--cycle;
  \draw[sdrAmber,line width=1.2pt] (1.5,0.45)--(1.5,1.35);
  \fill[sdrAmber!80] (1.8,0.45)--(1.88,1.4)--(2.62,1.4)--(2.65,0.45)--(2.68,1.4)--(3.42,1.4)--(3.5,0.45)--cycle;
  \fill[sdrGreen!60!white] (3.68,0.45) rectangle (3.98,0.85);
  \node[text=white,font=\tiny] at (0.75,0.25) {L+R}; \node[text=white,font=\tiny] at (1.5,0.25) {19};
  \node[text=white,font=\tiny] at (2.65,0.25) {L$-$R}; \node[text=white,font=\tiny] at (3.83,0.25) {RDS};}

\begin{document}
\SDRbanner

\begin{infostrip}
\begin{tabularx}{\linewidth}{@{}r@{\hspace{1em}}X@{}}
\textbf{دو جلسه} & جلسهٔ اول (بخش‌های ۱ و ۲؛ بخش ۳ اختیاری): دمدولاسیون همدوس و ناهمدوس و پیوند کابلی
% [orig] \textbf{دو جلسه} & جلسهٔ اول (بخش‌های ۱ تا ۳): دمدولاسیون همدوس و ناهمدوس و پیوند کابلی
% [orig] \textbf{دو جلسه} & جلسهٔ اول (بخش‌های ۱ تا ۵): دمدولاسیون همدوس و ناهمدوس و پیوند کابلی
\lr{NBFM}.
% جلسهٔ دوم (بخش‌های ۶ تا ۹):
% \lr{WFM}
% مونو و استریو و
% \lr{RDS}. \\
\\
\textbf{شرط ورود} & برای هر جلسه: گزارش پیش‌از‌جلسهٔ همان جلسه و کوئیز ورودی \\
\textbf{خروجی هر جلسه} & فایل‌های خام \IQ{} با نام استاندارد
(\code{lab02})،
فلوگراف‌ها، تصاویر نمایشگر، لاگ‌بوک، چک‌لیست راه‌اندازی امضاشده و برگهٔ تسویهٔ تجهیزات \\
\textbf{تازه در این آزمایش} & نخستین ارسال رادیویی درس: فقط کابلی و با تضعیف‌گر (فقط در بخش اختیاری ۳)
% [orig] \textbf{تازه در این آزمایش} & نخستین ارسال رادیویی درس: فقط کابلی و با تضعیف‌گر
\end{tabularx}
\end{infostrip}

%======================================================================
\isection{target}{اهداف آزمایش}
%======================================================================
\noindent\textbf{جلسهٔ اول.} تا پایان جلسه باید:
\begin{objectives}
\item یک ایستگاه
\lr{FM}
تجاری را با آشکارساز ناهمدوس و آشکارساز همدوس دریافت و کیفیت شنیداری و نسبت سیگنال به نویز خروجی آن دو را مقایسه کرده باشید؛
\item یک سیگنال
\lr{NBFM}
ساخت خودتان را از راه کابل و تضعیف‌گر فرستاده باشید (اختیاری)؛
% [orig] ساخت خودتان را از راه کابل و تضعیف‌گر فرستاده باشید؛
% [orig] ساخت خودتان را از راه کابل و تضعیف‌گر فرستاده و با
% \lr{SDR}
% دوم دریافت کرده باشید؛
% \item دادهٔ لازم برای رسم
% $\mathrm{SNR}$
% خروجی بر حسب
% $\mathrm{SNR}$
% ورودی و تعیین آستانهٔ
% \lr{FM}
% را ضبط کرده باشید؛
% \item در چالش تیمی، اثر انحراف فرکانس اختصاصی تیم را بر پهنای باند و کیفیت سنجیده باشید.
\end{objectives}
% \noindent\textbf{جلسهٔ دوم.} تا پایان جلسه باید:
% \begin{objectives}
% \item گیرندهٔ
% \lr{WFM}
% مونو ساخته و اجزای طیف باند پایهٔ استریو را با فرکانس اندازه‌گیری‌شده مشخص کرده باشید؛
% \item با بازیابی پایلوت و دو برابر کردن فرکانس آن، صدای استریو را بازسازی و جداسازی دو کانال را اندازه گرفته باشید؛
% \item \lr{RDS}
% را دیکد و نام ایستگاه
% \eng{PS}
% و شناسهٔ برنامه
% \eng{PI}
% را استخراج کرده باشید؛
% \item دادهٔ لازم برای تحلیل نرخ خطای دیکد
% \lr{RDS}
% در برابر سطح سیگنال را ضبط کرده باشید.
% \end{objectives}

%======================================================================
\isection{tools}{تجهیزات و نرم‌افزار مورد نیاز}
%======================================================================
\begin{table}[H]
	\centering\small
	\caption{تجهیزات هر تیم در آزمایش ۲}
	\label{tab:equip2}
	\begin{tabularx}{\linewidth}{|r|c|X|}
		\hline
		\hdrrow \hdr{قلم} & \hdr{تعداد} & \hdr{کاربرد} \\ \hline
		گیرندهٔ \lr{Pluto SDR} با کابل رابط \lr{USB} و آنتن & ۱ & گیرنده در هر دو جلسه \\ \hline
% [orig] 		گیرندهٔ \RTLSDR{} با کابل رابط \lr{USB} و آنتن \lr{VHF} & ۱ & گیرنده در هر دو جلسه \\ \hline
		\lr{SDR} فرستنده (\lr{ADALM-PLUTO}، یا \lr{HackRF}، یا \lr{USRP}) & ۱ & فرستندهٔ پیوند کابلی (بخش اختیاری ۳) \\ \hline
% [orig] 		\lr{SDR} فرستنده (\lr{ADALM-PLUTO}، یا \lr{HackRF}، یا \lr{USRP}) & ۱ & فرستندهٔ پیوند کابلی \\ \hline
		تضعیف‌گر ثابت \q{30}{dB} و کابل \lr{RF} & ۱ & اتصال فرستنده به گیرنده (بخش اختیاری ۳) \\ \hline
% [orig] 		تضعیف‌گر ثابت \q{30}{dB} و کابل \lr{RF} & ۱ & اتصال فرستنده به گیرنده (الزامی) \\ \hline
% 		بار ۵۰ اهمی & ۱ & بستن درگاه فرستنده هنگامی که به گیرنده وصل نیست \\ \hline
		هدفون & ۱ & شنیدن خروجی (به‌ویژه استریو) \\ \hline
		رایانه با \GR{} \lr{3.10}، \lr{Python} و \lr{Git} & ۱ & اجرای هم‌زمان فلوگراف فرستنده و گیرنده \\ \hline
% [orig] 		رایانه با \GR{} \lr{3.10}، ماژول \lr{gr-rds} (جلسهٔ دوم)، \lr{Python} و \lr{Git} & ۱ & اجرای هم‌زمان فلوگراف فرستنده و گیرنده \\ \hline
		فایل \code{sdrlab_iq_recorder.py} & — & بلوک ضبط \\ \hline
% [orig] 		فایل‌های \code{sdrlab_iq_recorder.py}، \code{sdrlab_rds_counter.py} و \code{lab02_stereo_rds_tx.grc} & — & بلوک ضبط، شمارندهٔ گروه \lr{RDS} و فرستندهٔ آمادهٔ جلسهٔ دوم \\ \hline
	\end{tabularx}
\end{table}

% \noindent
% فلوگراف فرستنده برای
% \lr{ADALM-PLUTO}
% نوشته شده است. برای سخت‌افزار دیگر فقط بلوک خروجی عوض می‌شود (جدول
% \ref{tab:txhw}).
% [INSTRUCTOR][HW] Labels checked against GNU Radio 3.10.9 block definitions.
% Maximum output powers are nominal data-sheet values; verify with your units.
% \begin{table}[H]
% 	\centering\small
% 	\caption{بلوک خروجی فرستنده برای سخت‌افزارهای مختلف (\GR{} \lr{3.10})}
% 	\label{tab:txhw}
% 	\begin{tabularx}{\linewidth}{|c|c|X|X|}
% 		\hline
% 		\hdrrow \hdr{سخت‌افزار} & \hdr{بلوک} & \hdr{فرکانس و توان} & \hdr{حد ایمن در این آزمایش} \\ \hline
% 		\lr{ADALM-PLUTO} & \blk{PlutoSDR Sink} & \prm{LO Frequency}؛ \prm{Attenuation TX1 (dB)} (بیشتر یعنی توان کمتر) & تضعیف دست‌کم \q{10}{dB}؛ بیشینهٔ خروجی نوعی \q{+7}{dBm}؛ بازهٔ پیش‌فرض \q{325}{MHz} تا \q{3.8}{GHz} \\ \hline
% 		\lr{HackRF One} & \blk{Soapy HackRF Sink} & \prm{Center Freq (Hz)}؛ \prm{VGA Gain (0dB - 47dB)} & \prm{Amp On (+14 dB)} همیشه خاموش؛ \prm{VGA Gain} حداکثر \val{20} \\ \hline
% 		\lr{USRP} & \blk{UHD: USRP Sink} & \prm{Ch0: Center Freq (Hz)}؛ \prm{Ch0: Gain Value} & بهره را از صفر آغاز کنید؛ بیشینهٔ توان را از برگهٔ مشخصات برد بخوانید \\ \hline
% 	\end{tabularx}
% \end{table}

%======================================================================
\isection{warn}{ایمنی و مقررات این آزمایش}
%======================================================================
قواعد همیشگی آزمایشگاه (سقف
\q{-10}{dBm}
در ورودی گیرنده، حساسیت به
\lr{DC}،
تخلیهٔ الکتریسیتهٔ ساکن) مانند آزمایش ۱ برقرار است. این آزمایش ارسال دارد و قواعد زیر به آن افزوده می‌شود. قواعد ارسال و چک‌لیست زیر فقط در صورت انجام بخش اختیاری ۳ لازم است.
% [orig] تخلیهٔ الکتریسیتهٔ ساکن) مانند آزمایش ۱ برقرار است. این آزمایش ارسال دارد و قواعد زیر به آن افزوده می‌شود.

\begin{safetybox}[قواعد ارسال (آیین‌نامه، بندهای \bnd{۸-۱} و \bnd{۸-۲})]
\begin{itemize}
\item هر ارسالی فقط از راه کابل و تضعیف‌گر انجام می‌شود. به درگاه فرستنده آنتن وصل نکنید.
\item خروجی فرستنده را هرگز مستقیم به ورودی گیرنده وصل نکنید؛ همیشه دست‌کم تضعیف‌گر
\q{30}{dB}
در مسیر باشد.
\item ارسال هوایی در باندهای دارای مجوز (از جمله باند پخش
\lr{FM})
تحت هیچ شرایطی مجاز نیست. فرکانس کار
$f_{\mathrm{tx}}$
را فقط از فهرست اعلامی مدرس بردارید.
\item پیش از هر ارسال، خالی‌بودن فرکانس هدف را با روش آزمایش ۱ بررسی کنید.
\item کابل
\lr{RF}
را فقط وقتی وصل یا جدا کنید که فلوگراف فرستنده متوقف است.
\end{itemize}
\end{safetybox}

\begin{cautionbox}[محاسبهٔ سطح پیش از اتصال]
توان در ورودی گیرنده برابر است با
\begin{equation}
	P_{\mathrm{rx}} = P_{\max} - A_{\mathrm{tx}} + 20\log_{10}(a_{\mathrm{tx}}) - A_{\mathrm{pad}} - L_{\mathrm{cable}}
	\label{eq:budget}
\end{equation}
که
$P_{\max}$
بیشینهٔ توان فرستنده،
$A_{\mathrm{tx}}$
تضعیف تنظیم‌شده در فرستنده،
$a_{\mathrm{tx}}$
دامنهٔ دیجیتال (میان صفر و یک)،
$A_{\mathrm{pad}}$
تضعیف‌گر ثابت و
$L_{\mathrm{cable}}$
تلفات کابل است. با مقادیر آغازین این دستور کار
($P_{\max} = +7\un{dBm}$، $A_{\mathrm{tx}} = 40\un{dB}$، $a_{\mathrm{tx}} = 0.5$، $A_{\mathrm{pad}} = 30\un{dB}$):
$P_{\mathrm{rx}} \approx 7 - 40 - 6 - 30 - 0.5 = -69.5\un{dBm}$.
حتی با کمترین تضعیف مجاز
($A_{\mathrm{tx}} = 10\un{dB}$)
توان
\q{-39.5}{dBm}
است و از سقف
\q{-10}{dBm}
فاصلهٔ کافی دارد. بدون تضعیف‌گر ثابت این فاصله از بین می‌رود.
\end{cautionbox}

\noindent
\textbf{چک‌لیست راه‌اندازی}
(اگر بخش ۳ انجام می‌شود، پیش از نخستین ارسال کامل و به امضای دستیار آموزشی برسد):
% [orig] (در هر دو جلسه پیش از نخستین ارسال کامل و به امضای دستیار آموزشی برسد):
\begin{checklist}
\item مسیر «درگاه فرستنده، تضعیف‌گر
\q{30}{dB}،
کابل، ورودی گیرنده» بسته شده و هیچ آنتنی به فرستنده وصل نیست.
\item فرکانس کار تیم از فهرست مدرس انتخاب و خالی‌بودن آن با گیرنده بررسی و در لاگ‌بوک ثبت شده است.
\item در فلوگراف فرستنده، تضعیف آغازین
\q{40}{dB} (بدون تضعیف‌گر خارجی: \q{70}{dB})
% [orig] \q{40}{dB}
و دامنهٔ دیجیتال حداکثر
\val{0.5}
است.
\item هر سه عضو می‌توانند رابطهٔ
\eqref{eq:budget}
را برای تنظیمات فعلی حساب کنند.
\item در گیرنده \prm{Gain Mode (RX1)} روی \val{manual} (بدون تقویت خودکار) است.
% [orig] \item در گیرنده
% [orig] \prm{Bias Tee Power}
% [orig] و
% [orig] \prm{AGC}
% [orig] خاموش است.
\end{checklist}
\noindent
{\small فرکانس کار تیم: $f_{\mathrm{tx}} = $ \blank[2.6cm] \lr{MHz} \hfill امضای دستیار آموزشی: \blank[3.4cm] \hfill ساعت: \blank[1.6cm]}

%======================================================================
% \isection{clock}{نقشهٔ زمانی دو جلسه}
%======================================================================
% \begin{table}[H]
% 	\centering\small
% 	\caption{زمان‌بندی دو جلسه (مطابق الگوی آیین‌نامه)}
% 	\label{tab:timeline2}
% 	\begin{tabularx}{\linewidth}{|c|X|X|}
% 		\hline
% 		\hdrrow \hdr{دقیقه} & \hdr{جلسهٔ اول} & \hdr{جلسهٔ دوم} \\ \hline
% 		۰ تا ۱۵ & کوئیز ورودی و کنترل گزارش پیش‌از‌جلسه & کوئیز ورودی و کنترل گزارش پیش‌از‌جلسه \\ \hline
% 		۱۵ تا ۳۰ & توضیح مدرس، مرور ایمنی ارسال، چک‌لیست راه‌اندازی & توضیح مدرس، مرور ایمنی، چک‌لیست راه‌اندازی \\ \hline
% 		۳۰ تا ۵۵ & بخش \ref{part:rx1}: گیرنده با آشکارساز ناهمدوس & بخش \ref{part:mono}: \lr{WFM} مونو و طیف باند پایه \\ \hline
% 		۵۵ تا ۷۵ & بخش \ref{part:pll}: آشکارساز همدوس و مقایسه & بخش \ref{part:stereo}: بازسازی استریو (تا دقیقهٔ ۹۰) \\ \hline
% 		۷۵ تا ۱۰۰ & بخش \ref{part:link}: پیوند کابلی \lr{NBFM} & بخش \ref{part:rds}: دیکد \lr{RDS} (دقیقهٔ ۹۰ تا ۱۲۰) \\ \hline
% 		۱۰۰ تا ۱۲۰ & بخش \ref{part:thr}: $\mathrm{SNR}$ خروجی و آستانه & \\ \hline
% 		۱۲۰ تا ۱۵۰ & بخش \ref{part:dev}: چالش تیمی (انحراف اختصاصی) & بخش \ref{part:rdslevel}: چالش تیمی (\lr{RDS} در برابر سطح) \\ \hline
% 		۱۵۰ تا ۱۸۰ & پرسش شفاهی، جمع‌بندی، تحویل تجهیزات & پرسش شفاهی، جمع‌بندی، تحویل تجهیزات \\ \hline
% 	\end{tabularx}
% \end{table}

%======================================================================
\isection{book}{مرور کوتاه تئوری}
%======================================================================
جزئیات و جدول «ایده‌آل نظری در برابر پیاده‌سازی عملی» در راهنمای پیش‌از‌جلسه آمده است. با نرخ
$f_q$
در ورودی آشکارساز، بیشینهٔ انحراف
$\Delta f$
و پهنای باند صوتی
$W$:
\begin{equation}
	B_T \approx 2(\Delta f + W), \qquad
	G_{\mathrm{det}} = \frac{f_q}{2\pi \Delta f}, \qquad
	k_f = \frac{2\pi \Delta f}{f_q}, \qquad
	B_L \approx 1.06 \times (\text{\lr{Loop Bandwidth}}) \times f_q
	\label{eq:recap2}
\end{equation}
\begin{equation}
	\mathrm{CNR} = \frac{P_{\mathrm{on}} - P_{\mathrm{off}}}{P_{\mathrm{off}}}, \qquad
	\mathrm{SNR}_o \approx \frac{3}{2}\left(\frac{\Delta f}{W}\right)^{2} \frac{B_T}{W}\, \mathrm{CNR}
	\quad \text{(\fa{بالاتر از آستانه})}, \qquad
	\mathrm{SINAD} = \frac{P_{S+N+D}}{P_{N+D}}
	\label{eq:snr2}
\end{equation}
که
$P_{\mathrm{on}}$
و
$P_{\mathrm{off}}$
توان خطی کانال با فرستندهٔ روشن و خاموش‌اند (نمایشگر آن‌ها را بر حسب
\dBFS{}
می‌دهد؛ پیش از تفریق به مقیاس خطی ببرید). حامل تک‌تن در
$\Delta f = 2.405\, f_m$
و
$5.520\, f_m$
صفر می‌شود.

%======================================================================
\isection{flow}{معماری سیستم و فلوگراف‌ها}
%======================================================================
\subsection{چیدمان سخت‌افزار و برنامهٔ نرخ}
در این آزمایش دو چیدمان به کار می‌رود (شکل
\ref{fig:setup2}):
دریافت از آنتن برای ایستگاه‌های واقعی، و پیوند کابلی میان فرستنده و گیرنده (فقط برای بخش اختیاری ۳). هر دو
% [orig] دریافت از آنتن برای ایستگاه‌های واقعی، و پیوند کابلی میان فرستنده و گیرنده. هر دو
\lr{SDR}
به یک رایانه وصل‌اند و دو فلوگراف هم‌زمان اجرا می‌شوند. همهٔ فلوگراف‌ها یک برنامهٔ نرخ دارند:
\q{2.4}{MS/s}
در سخت‌افزار،
\q{240}{kHz}
در آشکارساز و
\q{48}{kHz}
در صوت.

\begin{figure}[H]
	\centering
	\begin{tikzpicture}
		\node[font=\footnotesize\bfseries,text=sdrDark] at (-1.6,1.3) {\fa{(الف)}};
		\node[sdrblkhw,minimum width=2.3cm] (ant) at (1.0,1.3) {\fa{آنتن}};
% [orig] 		\node[sdrblkhw,minimum width=2.3cm] (ant) at (1.0,1.3) {\fa{آنتن \lr{VHF}}};
		\node[sdrblkhl,minimum width=2.2cm] (rx1) at (5.0,1.3) {Pluto SDR};
% [orig] 		\node[sdrblkhl,minimum width=2.2cm] (rx1) at (5.0,1.3) {RTL-SDR};
		\node[sdrblkpc,minimum width=2.6cm] (pc1) at (9.6,1.3) {\fa{رایانه: \lr{lab02\_fm\_rx}}};
		\draw[sdrarr] (ant)--(rx1); \draw[cplx] (rx1)--node[above,font=\scriptsize,text=black]{USB}(pc1);
		\node[font=\footnotesize\bfseries,text=sdrDark] at (-1.6,-0.6) {\fa{(ب)}};
		\node[sdrblkhw,minimum width=2.3cm] (tx) at (1.0,-0.6) {\fa{فرستنده}\\ (PlutoSDR)};
		\node[sdrblk,minimum width=1.9cm] (pad) at (4.1,-0.6) {\fa{تضعیف‌گر}\\ 30 dB};
		\node[sdrblkhl,minimum width=2.0cm] (rx2) at (7.1,-0.6) {Pluto SDR};
% [orig] 		\node[sdrblkhl,minimum width=2.0cm] (rx2) at (7.1,-0.6) {RTL-SDR};
		\node[sdrblkpc,minimum width=2.6cm] (pc2) at (11.0,-0.6) {\fa{رایانه: \lr{TX + RX}}};
		\draw[sdrarr] (tx)--node[above,font=\scriptsize]{\fa{کابل}}(pad); \draw[sdrarr] (pad)--(rx2);
		\draw[cplx] (rx2)--node[above,font=\scriptsize,text=black]{USB}(pc2);
		\draw[cplx] (pc2.south) -- ++(0,-0.55) -| node[pos=0.25,below,font=\scriptsize,text=black]{USB}(tx.south);
	\end{tikzpicture}
	\caption{دو چیدمان آزمایش: (الف) دریافت از آنتن، (ب) پیوند کابلی با تضعیف‌گر (اختیاری)}
% [orig] 	\caption{دو چیدمان آزمایش: (الف) دریافت از آنتن، (ب) پیوند کابلی با تضعیف‌گر}
	\label{fig:setup2}
\end{figure}

\begin{notebox}[تنظیم با آفست]
گیرنده را
\q{250}{kHz}
پایین‌تر از سیگنال تنظیم می‌کنیم
(\code{f_off})
و فیلتر
\blk{Frequency Xlating FIR Filter}
سیگنال را به مرکز می‌آورد. به این ترتیب قلهٔ
\lr{DC}
گیرنده بیرون از کانال می‌ماند. فرستنده نیز سیگنال را
\q{250}{kHz}
بالاتر از فرکانس نوسان‌ساز خود می‌سازد تا نشت حامل فرستنده روی سیگنال نیفتد.
\end{notebox}

\subsection{گیرندهٔ جلسهٔ اول: \lr{lab02\_fm\_rx}}
شکل
\ref{fig:fg-rx}
مسیر سیگنال و جدول‌های
\ref{tab:params-rx}
و
\ref{tab:params-rx2}
پارامترها را می‌دهند. دو آشکارساز به‌موازات هم کار می‌کنند و
\blk{Selector}
یکی را به زنجیرهٔ صوت می‌دهد. سه اندازه‌گیر نیز در فلوگراف هست: توان کانال (برای
$\mathrm{CNR}$)،
خطای فرکانس حامل (میانگین خروجی تمایزگر ضرب در
$\Delta f$)
و
\lr{SINAD}.

\begin{figure}[H]
	\centering
	\begin{tikzpicture}
		\node[grcsrc] (src) at (0,0) {\grcname{PlutoSDR Source}\\ \grcp{Sample Rate}{2.4e6}\\ \grcp{LO Frequency}{fc - f\_off}};
% [orig] 		\node[grcsrc] (src) at (0,0) {\grcname{PlutoSDR Source}\\ \grcp{Sample Rate}{2.4e6}\\ \grcp{Center Freq}{fc - f\_off}};
% [orig] 		\node[grcsrc] (src) at (0,0) {\grcname{Soapy RTLSDR Source}\\ \grcp{Sample Rate}{2.4e6}\\ \grcp{Center Freq}{fc - f\_off}};
		\node[grc,anchor=west] (xl) at (3.15,0) {\grcname{Freq. Xlating FIR}\\ \grcp{Decimation}{10}\\ \grcp{Center}{f\_off + f\_fine}};
		\node[grc,anchor=west] (cf) at (6.6,0) {\grcname{Low Pass Filter}\\ \grcp{Cutoff}{chan\_bw/2}};
		\node[grc,anchor=west] (qd) at (9.7,0.75) {\grcname{Quadrature Demod}};
		\node[grc,anchor=west] (pl) at (9.7,-0.75) {\grcname{PLL Freq. Detector}\\ $\times$ \texttt{Gain}};
		\node[grc,anchor=west] (sel) at (13.1,0) {\grcname{Selector}\\ \grcp{in}{det}};
		\node[grc,anchor=west] (af) at (6.6,-3.25) {\grcname{Low Pass Filter}\\ \grcp{Decimation}{5}\\ \grcp{Cutoff}{audio\_bw}};
		\node[grc,anchor=west] (dc) at (9.7,-3.25) {\grcname{DC Blocker}};
		\node[grcsink,anchor=west] (au) at (12.1,-2.75) {\grcname{FM Deemphasis}\\ $\rightarrow$ \grcname{Audio Sink}};
		\node[grcsink,anchor=west] (sn) at (12.1,-3.95) {\grcname{SINAD meter}\\ RMS, Band Reject, Log10};
		\node[grcsink,anchor=north] (pm) at (7.9,-0.95) {\grcname{RMS $\to$ Log10}: \texttt{P\_ch}};
		\node[grcfile,anchor=north] (rec) at (4.2,-0.95) {\grcname{IQ Recorder}\\ \grcp{fs}{240e3}};
		\draw[cplx] (src)--(xl); \draw[cplx] (xl)--(cf);
		\draw[cplx] (cf.east) -- ++(0.25,0) |- (qd.west); \draw[cplx] (cf.east) -- ++(0.25,0) |- (pl.west);
		\draw[flt] (qd.east) -- ++(0.3,0) |- ([yshift=3pt]sel.west); \draw[flt] (pl.east) -- ++(0.3,0) |- ([yshift=-3pt]sel.west);
		\draw[flt] (sel.south) |- ++(-7.95,-1.78) |- (af.west);
		\draw[flt] (af)--(dc); \draw[flt] (dc.east) -- ++(0.2,0) |- (au.west); \draw[flt] (dc.east) -- ++(0.2,0) |- (sn.west);
		\draw[cplx] (cf.south)--(pm.north -| cf.south); \draw[cplx] (xl.south)--(rec.north -| xl.south);
	\end{tikzpicture}
	\caption{طرح‌وارهٔ گیرندهٔ جلسهٔ اول (رسم‌شده؛ تصویر واقعی \GRC{} نیست). آبی: جریان مختلط، نارنجی: جریان حقیقی}
	\label{fig:fg-rx}
\end{figure}
% \screenshotslot{labs/figs/lab02_fm_rx.png}{تصویر فلوگراف \lr{lab02\_fm\_rx} در \GRC}

\begin{table}[H]
	\centering\small
	\caption{پارامترهای گیرندهٔ \lr{lab02\_fm\_rx} (\GR{} \lr{3.10})}
	\label{tab:params-rx}
	\begin{tabularx}{\linewidth}{|C{4.4cm}|Z|}
		\hline
		\hdrrow \hdr{بلوک (\lr{Id})} & \hdr{پارامترها} \\ \hline
		\blk{Variable} & \pv{samp\_rate}{2.4e6}؛ \pv{quad\_rate}{240e3}؛ \pv{audio\_rate}{48e3}؛ \pv{f\_off}{250e3}؛ \pv{chan\_bw}{min(2*(dev + audio\_bw), 200e3)}؛ \pv{team}{"Team-07"} \\ \hline
		\blk{Import} & \pv{Import}{import math} \\ \hline
		\blk{QT GUI Range} & \code{fc}: فرکانس سیگنال (\val{98.0e6}، گام \val{100e3})؛ \code{f_fine}: تنظیم ریز (\val{0}، از \val{-30e3} تا \val{30e3}، گام \val{10})؛ \code{rf_gain} (\val{20.7})؛ \code{dev} (\val{75e3}، از \val{1e3} تا \val{75e3})؛ \code{f_m} (\val{1e3})؛ \code{pll_bw} (\val{0.3}، از \val{0.005} تا \val{0.6})؛ \code{volume} (\val{0.3}) \\ \hline
		\blk{QT GUI Chooser} & \code{det}: \val{0} (\lr{Quadrature Demod}) یا \val{1} (\lr{PLL})؛ \code{audio_bw}: \val{15e3} یا \val{3e3}؛ \code{tau}: \val{50e-6} یا \val{75e-6} \\ \hline
		\blk{PlutoSDR Source} & \pv{Sample Rate}{samp\_rate}؛ \pv{LO Frequency}{int(fc - f\_off)}؛ \pv{Gain Mode (RX1)}{manual}؛ \pv{Manual Gain (RX1)(dB)}{rf\_gain} \\ \hline
% [orig] 		\blk{Soapy RTLSDR Source} & \pv{Sample Rate}{samp\_rate}؛ \pv{Center Freq (Hz)}{fc - f\_off}؛ \pv{AGC}{False}؛ \pv{RF Gain}{rf\_gain}؛ \pv{Bias Tee Power}{False} \\ \hline
		\blk{Frequency Xlating}\newline\blk{FIR Filter} & \pv{Type}{Complex->Complex (Complex Taps)}؛ \pv{Decimation}{10}؛ \pv{Taps}{firdes.low\_pass(1, samp\_rate, 100e3, 20e3)}؛ \pv{Center Frequency}{f\_off + f\_fine}؛ \pv{Sample Rate}{samp\_rate} \\ \hline
		\blk{Low Pass Filter} (کانال) & \pv{FIR Type}{Complex->Complex (Decimating)}؛ \pv{Decimation}{1}؛ \pv{Sample Rate}{quad\_rate}؛ \pv{Cutoff Freq}{chan\_bw/2}؛ \pv{Transition Width}{chan\_bw/8} \\ \hline
		\blk{Quadrature Demod} & \pv{Gain}{quad\_rate/(2*math.pi*dev)} \\ \hline
		\blk{PLL Frequency Detector} & \pv{Loop Bandwidth}{pll\_bw}؛ \pv{Max Phase/sample}{min(2*math.pi*1.5*dev/quad\_rate, 3.0)}؛ \prm{Min Phase/sample}: همان با علامت منفی \\ \hline
		\blk{Multiply Const} (پس از \lr{PLL}) & \pv{IO Type}{float}؛ \pv{Constant}{quad\_rate/(2*math.pi*dev)} \\ \hline
		\blk{Selector} & \pv{Type}{float}؛ \pv{Number of Inputs}{2}؛ \pv{Number of Outputs}{1}؛ \pv{Input Index}{det} \\ \hline
		\blk{Low Pass Filter} (صوت) & \pv{FIR Type}{Float->Float (Decimating)}؛ \pv{Decimation}{5}؛ \pv{Sample Rate}{quad\_rate}؛ \pv{Cutoff Freq}{audio\_bw}؛ \pv{Transition Width}{0.25*audio\_bw + 250} \\ \hline
		\blk{DC Blocker} & \pv{Type}{Float->Float}؛ \pv{Length}{512} \\ \hline
		\blk{FM Deemphasis}، \blk{Multiply Const}، \blk{Audio Sink} & \pv{Sample Rate}{audio\_rate}؛ \pv{Tau}{tau}؛ سپس ضرب در \code{volume} و \blk{Audio Sink} با \pv{Sample Rate}{48 kHz} \\ \hline
	\end{tabularx}
\end{table}
\begin{table}[H]
	\centering\small
	\caption{اندازه‌گیرها، نمایشگرها و ضبط در گیرندهٔ \lr{lab02\_fm\_rx}}
	\label{tab:params-rx2}
	\begin{tabularx}{\linewidth}{|C{4.4cm}|Z|}
		\hline
		\hdrrow \hdr{بخش} & \hdr{بلوک‌ها و پارامترها} \\ \hline
		توان کانال & خروجی فیلتر کانال به \blk{RMS} (\pv{Input Type}{complex}؛ \pv{Alpha}{2e-5})، \blk{Log10} (\pv{n}{20}) و \blk{QT GUI Number Sink} \\ \hline
		خطای فرکانس & خروجی \blk{Quadrature Demod} به \blk{Single Pole IIR Filter} (\pv{Alpha}{2e-5})، \blk{Multiply Const} (\pv{Constant}{dev}) و \blk{QT GUI Number Sink} \\ \hline
		\lr{SINAD} & خروجی \blk{DC Blocker} از دو مسیر به \blk{Subtract} می‌رود: (۱) \blk{RMS} (\pv{Alpha}{2e-4}) و \blk{Log10} (\pv{n}{20})؛ (۲) \blk{Band Reject Filter} (\pv{Low Cutoff Freq}{f\_m - 300}؛ \pv{High Cutoff Freq}{f\_m + 300}؛ \pv{Transition Width}{200}) و سپس همان \blk{RMS} و \blk{Log10}. خروجی \blk{Subtract} به \blk{QT GUI Number Sink} \\ \hline
		نمایشگرها & \blk{QT GUI Frequency Sink} مختلط روی خروجی فیلتر کانال (\pv{Bandwidth (Hz)}{quad\_rate}؛ \pv{Center Frequency (Hz)}{fc})؛ \blk{QT GUI Frequency Sink} حقیقی با دو ورودی روی خروجی دو آشکارساز (\pv{FFT Size}{4096})؛ \blk{QT GUI Frequency Sink} حقیقی روی صوت (\pv{Bandwidth (Hz)}{audio\_rate}) \\ \hline
		\blk{Python Block} (ضبط) & کد آزمایش ۱؛ ورودی از خروجی \blk{Frequency Xlating FIR Filter}؛ \pv{Lab}{2}؛ \pv{Fc}{fc}؛ \pv{Fs}{quad\_rate}؛ \pv{Gain}{rf\_gain}؛ \pv{Record}{rec}؛ \pv{Max Seconds}{30} \\ \hline
	\end{tabularx}
\end{table}

\begin{notebox}[ضبط در نرخ \q{240}{kHz}]
در این آزمایش جریان
\IQ{}
پس از انتخاب کانال (نرخ
\q{240}{kHz})
ضبط می‌شود، نه جریان
\q{2.4}{MS/s}.
حجم فایل ده برابر کمتر است (حدود
\q{1.9}{MB}
در ثانیه) و کل سیگنال
\lr{FM}
را در بر دارد. در نام فایل،
\code{fs}
برابر
\val{240000}
و
\code{fc}
فرکانس خود سیگنال است.
\end{notebox}

\subsection{فرستندهٔ کابلی: \lr{lab02\_fm\_tx} (بخش اختیاری ۳)}
% [orig] \subsection{فرستندهٔ کابلی: \lr{lab02\_fm\_tx}}
% \begin{center}
% \lr{\small\sffamily Signal Source \then\ Frequency Mod \then\ Rational Resampler ($\times 10$) \then\ Rotator \then\ Multiply Const \then\ PlutoSDR Sink}
% \end{center}

\begin{table}[H]
	\centering\small
	\caption{پارامترهای فرستندهٔ \lr{lab02\_fm\_tx} (اختیاری)}
% [orig] 	\caption{پارامترهای فرستندهٔ \lr{lab02\_fm\_tx}}
	\label{tab:params-tx}
	\begin{tabularx}{\linewidth}{|C{4.4cm}|Z|}
		\hline
		\hdrrow \hdr{بلوک (\lr{Id})} & \hdr{پارامترها} \\ \hline
		\blk{Variable} & \pv{samp\_rate}{2.4e6}؛ \pv{quad\_rate}{240e3}؛ \pv{f\_off}{250e3} \\ \hline
		\blk{QT GUI Range} & \code{f_tx}: فرکانس کار تیم (پیش‌فرض \val{434.0e6})؛ \code{dev} (\val{5e3}، از \val{500} تا \val{75e3}، گام \val{50})؛ \code{f_m} (\val{1e3})؛ \code{a_m} (\val{1.0})؛ \code{tx_atten} (\val{40}، از \val{10} تا \val{89}، گام \val{1})؛ \code{tx_amp} (\val{0.5}، از \val{0} تا \val{0.8}) \\ \hline
		\blk{Signal Source} & \pv{Output Type}{Float}؛ \pv{Sample Rate}{quad\_rate}؛ \pv{Waveform}{Cosine}؛ \pv{Frequency}{f\_m}؛ \pv{Amplitude}{a\_m} \\ \hline
		\blk{Frequency Mod} & \pv{Sensitivity}{2*math.pi*dev/quad\_rate} \\ \hline
		\blk{Rational Resampler} & \pv{Type}{Complex->Complex (Complex Taps)}؛ \pv{Interpolation}{10}؛ \pv{Decimation}{1} \\ \hline
		\blk{Rotator} & \pv{Phase Increment}{2*math.pi*f\_off/samp\_rate} \\ \hline
		\blk{Multiply Const} & \pv{IO Type}{complex}؛ \pv{Constant}{tx\_amp} \\ \hline
		\blk{PlutoSDR Sink} & \pv{LO Frequency}{int(f\_tx - f\_off)}؛ \pv{Sample Rate}{int(samp\_rate)}؛ \pv{Attenuation TX1 (dB)}{tx\_atten}؛ \pv{RF Bandwidth (Hz)}{2000000} \\ \hline
		\blk{QT GUI Frequency Sink} & مختلط، روی خروجی \blk{Frequency Mod}؛ \pv{FFT Size}{4096}؛ \pv{Bandwidth (Hz)}{quad\_rate} \\ \hline
	\end{tabularx}
\end{table}
\noindent
تضعیف فرستنده
(\code{tx_atten})
هنگام اجرا تغییر می‌کند و با آن سطح دریافتی را پله‌پله کم خواهید کرد. برای فرستادن صدا به‌جای تن،
\blk{Signal Source}
را با
\blk{Wav File Source}
(\q{48}{kHz}،
تک‌کاناله) و یک
\blk{Rational Resampler}
با
\pv{Interpolation}{5}
جایگزین کنید.

%======================================================================
\isection{steps}{شرح آزمایش: جلسهٔ اول}
%======================================================================

\labpart[]{گیرندهٔ \lr{FM} با آشکارساز ناهمدوس}\label{part:rx1}
% [orig] \labpart[۲۵ دقیقه]{گیرندهٔ \lr{FM} با آشکارساز ناهمدوس}\label{part:rx1}
\begin{steps}
\item آنتن را به گیرنده ببندید (چیدمان الف) و فلوگراف
\code{lab02_fm_rx.grc}
را اجرا کنید:
\pv{dev}{75e3}،
پهنای صوتی
\q{15}{kHz}،
آشکارساز
\lr{Quadrature Demod}
و
\code{fc}
برابر فرکانس یک ایستگاه قوی از فهرست آزمایش ۱.
\item بهره را مانند آزمایش ۱ تنظیم کنید (بدون اشباع) و صدا را با هدفون بشنوید. ثابت زمانی واتأکید محل را انتخاب کنید و با مقدار دیگر مقایسه کنید.
\item در نمایشگر «خروجی آشکارسازها» (تا
\q{120}{kHz})
آنچه را بالاتر از باند صوتی می‌بینید یادداشت کنید.
% [orig] آنچه را بالاتر از باند صوتی می‌بینید یادداشت کنید؛ در جلسهٔ دوم به آن بازمی‌گردید.
\item \textbf{خطای فرکانس.} عدد «خطای فرکانس حامل» را بخوانید. با
\code{f_fine}
آن را به نزدیک صفر برسانید و مقدار
\code{f_fine}
را بنویسید. خطای گیرنده را بر حسب
\lr{ppm}
حساب کنید.
\item \textbf{توان کانال.} «توان کانال» را در بهرهٔ کاری بخوانید؛ سپس
\code{fc}
را روی یک کانال خالی مجاور بگذارید و دوباره بخوانید. از این دو عدد
$\mathrm{CNR}$
ایستگاه را در
\q{200}{kHz}
حساب کنید.
\item یک ضبط ده‌ثانیه‌ای از ایستگاه بگیرید.
\end{steps}

\begin{expectbox}
صدای روشن و بدون خش. خطای فرکانس گیرنده معمولاً از چند صد هرتز تا چند کیلوهرتز است و پس از تنظیم
\code{f_fine}
نزدیک صفر می‌شود. در طیف خروجی آشکارساز، بالاتر از صوت یک خط باریک در
\q{19}{kHz}
دیده می‌شود (اگر ایستگاه استریو باشد).
\end{expectbox}

\begin{question}
اگر در گیرنده
\code{dev}
را
\val{37.5e3}
بگذارید، سطح صدا چه تغییری می‌کند و چرا؟ چرا طیف
\lr{RF}
کانال با این تغییر عوض نمی‌شود؟
\end{question}
\begin{question}
چرا خطای فرکانس حامل در خروجی تمایزگر به‌صورت مقدار ثابت دیده می‌شود؟ خطای اندازه‌گیری‌شده با مقداری که در آزمایش ۱ برای همین گیرنده دیدید سازگار است؟
\end{question}

\labpart[]{آشکارساز همدوس و مقایسهٔ دو آشکارساز}\label{part:pll}
% [orig] \labpart[۲۰ دقیقه]{آشکارساز همدوس و مقایسهٔ دو آشکارساز}\label{part:pll}
\begin{steps}
\item آشکارساز را روی
\lr{PLL}
بگذارید
(\pv{pll\_bw}{0.3})
و بشنوید. میان دو آشکارساز چند بار جابه‌جا شوید.
\item \code{pll_bw}
را از
\val{0.5}
پله‌پله تا
\val{0.02}
کم کنید. در هر پله جدول
\ref{tab:pllbw}
را پر کنید: کیفیت صدا، و سطح خط
\q{19}{kHz}
در خروجی
\lr{PLL}
نسبت به خروجی تمایزگر (از نمایشگر دو منحنی).
\item \textbf{سیگنال ضعیف.}
\code{pll_bw}
را به کوچک‌ترین مقداری که صدا هنوز سالم است برگردانید. بهره را کم کنید تا نویز شنیده شود. در سه سطح سیگنال، برای هر آشکارساز سطح نویز را در ناحیهٔ خالی حوالی
\q{17}{kHz}
و سطح خط پایلوت را بخوانید (جدول
\ref{tab:weak}).
در هر سطح یک ضبط ده‌ثانیه‌ای بگیرید.
\end{steps}

\begin{table}[H]
	\centering\small
	\caption{اثر پهنای حلقه بر خروجی آشکارساز \lr{PLL} (ایستگاه: \blank[1.8cm] \lr{MHz})}
	\label{tab:pllbw}
	\begin{tabularx}{\linewidth}{|c|Y|Y|Y|Y|}
		\hline
		\hdrrow \hdr{\lr{pll\_bw}} & \hdr{$B_L$ (\lr{kHz})} & \hdr{صدا سالم است؟} & \hdr{خط \q{19}{kHz}: \lr{PLL} منهای تمایزگر (\lr{dB})} & \hdr{مؤلفه‌های بالای \q{23}{kHz} دیده می‌شوند؟} \\ \hline
		$0.5$ & \rowstrut & & & \\ \hline
		$0.3$ & \rowstrut & & & \\ \hline
		$0.15$ & \rowstrut & & & \\ \hline
		$0.08$ & \rowstrut & & & \\ \hline
		$0.04$ & \rowstrut & & & \\ \hline
		$0.02$ & \rowstrut & & & \\ \hline
	\end{tabularx}
\end{table}

\begin{table}[H]
	\centering\small
	\caption{مقایسهٔ دو آشکارساز در سطوح مختلف سیگنال (\lr{pll\_bw} = \blank[1.2cm])}
	\label{tab:weak}
	\begin{tabularx}{\linewidth}{|c|c|Y|Y|Y|Y|Y|}
		\hline
		\hdrrow & & \multicolumn{2}{c|}{\hdr{تمایزگر}} & \multicolumn{2}{c|}{\hdr{\lr{PLL}}} & \\ \hline
		\subhdrrow \textbf{بهره (\lr{dB})} & \textbf{توان کانال (\lr{dBFS})} & پایلوت (\lr{dB}) & نویز \q{17}{kHz} (\lr{dB}) & پایلوت (\lr{dB}) & نویز \q{17}{kHz} (\lr{dB}) & کدام بهتر شنیده می‌شود؟ \\ \hline
		\rowstrut & & & & & & \\ \hline
		\rowstrut & & & & & & \\ \hline
		\rowstrut & & & & & & \\ \hline
	\end{tabularx}
\end{table}

\begin{expectbox}
با حلقهٔ پهن
(\val{0.3}
و بیشتر) دو آشکارساز تقریباً یکسان شنیده می‌شوند. با باریک‌شدن حلقه ابتدا مؤلفه‌های بالای باند پایه
(\q{38}{}
و
\q{57}{kHz})
و سپس پایلوت در خروجی
\lr{PLL}
تضعیف می‌شوند؛ در شبیه‌سازی، حلقهٔ باریک‌تر از حدود
\val{0.1}
دیگر انحراف
$\pm 75\un{kHz}$
را دنبال نمی‌کند و صدا خراب می‌شود. در سیگنال ضعیف، حلقهٔ باریک‌تر نویز کمتری در خروجی دارد.
\end{expectbox}

% [orig block follows]
% \begin{question}
% چرا خروجی
% \lr{PLL}
% با حلقهٔ باریک مؤلفه‌های
% \q{38}{}
% و
% \q{57}{kHz}
% را از دست می‌دهد؟ نتیجه بگیرید برای دیکد استریو و
% \lr{RDS}
% در جلسهٔ دوم کدام آشکارساز مناسب‌تر است.
% \end{question}
\begin{question}
چرا خروجی
\lr{PLL}
با حلقهٔ باریک مؤلفه‌های
\q{38}{}
و
\q{57}{kHz}
را از دست می‌دهد؟
% [orig] \lr{RDS}
% [orig] در جلسهٔ دوم کدام آشکارساز مناسب‌تر است.
\end{question}
\begin{question}
در سیگنال ضعیف، نسبت «پایلوت به نویز» برای کدام آشکارساز بهتر بود؟ این نتیجه چه ارتباطی با نسبت پهنای حلقه به پهنای کانال دارد؟
\end{question}

\labpart[]{پیوند کابلی \lr{NBFM} ساخت خودتان (اختیاری)}\label{part:link}
% [orig] \labpart[]{پیوند کابلی \lr{NBFM} ساخت خودتان}\label{part:link}
% [orig] \labpart[۲۵ دقیقه]{پیوند کابلی \lr{NBFM} ساخت خودتان}\label{part:link}
\noindent\textbf{این بخش اختیاری است و در صورت وجود وقت و تجهیزات انجام می‌شود.}
\begin{steps}
\item \textbf{بررسی فرکانس.} پیش از هر ارسال، گیرنده را با آنتن روی
$f_{\mathrm{tx}}$
بگذارید
(\pv{dev}{5e3}،
پهنای صوتی
\q{3}{kHz})
و مطمئن شوید کانال خالی است. نتیجه را در لاگ‌بوک بنویسید.
\item \textbf{اتصال.} در حالی که فلوگراف فرستنده متوقف است، چیدمان (ب) را ببندید و چک‌لیست راه‌اندازی را به امضا برسانید.
\item «توان کانال» را با فرستندهٔ خاموش بخوانید
($P_{\mathrm{off}}$).
\item فلوگراف
\code{lab02_fm_tx.grc}
را اجرا کنید:
\pv{dev}{5e3}،
\pv{f\_m}{1e3}،
\pv{tx\_atten}{40}.
در گیرنده، خطای فرکانس را با
\code{f_fine}
صفر کنید (اگر سیگنال در کانال دیده نمی‌شود، موقتاً
\code{dev}
گیرنده را بزرگ کنید تا فیلتر کانال پهن شود).
\item طیف کانال را با
\pv{FFT Size}{4096}
ببینید: حامل و خطوط جانبی به فاصلهٔ
\q{1}{kHz}.
«توان کانال»
($P_{\mathrm{on}}$)
و
\lr{SINAD}
را بخوانید و تن را بشنوید.
\item \textbf{صفرشدن حامل.} با
\pv{f\_m}{1e3}،
\code{dev}
فرستنده را از
\val{1e3}
آهسته بالا ببرید و دو مقداری را که خط حامل در آن‌ها ناپدید می‌شود بیابید (جدول
\ref{tab:null}).
\item \textbf{پهنای باند.} برای دو حالت جدول
\ref{tab:null}
فاصلهٔ بیرونی‌ترین خطوط جانبی‌ای را که کمتر از
\q{20}{dB}
پایین‌تر از بزرگ‌ترین خط‌اند بخوانید و برای هر حالت یک ضبط پنج‌ثانیه‌ای بگیرید.
\end{steps}

\begin{table}[H]
	\centering\small
	\caption{کالیبراسیون انحراف با صفرشدن حامل و اندازه‌گیری پهنای باند ($f_{\mathrm{tx}} = $ \blank[1.8cm] \lr{MHz})}
	\label{tab:null}
	\begin{tabularx}{\linewidth}{|X|c|c|c|}
		\hline
		\hdrrow \hdr{کمیت} & \hdr{نظری} & \hdr{اندازه‌گیری} & \hdr{اختلاف} \\ \hline
		نخستین انحرافی که حامل صفر می‌شود، $f_m = 1\un{kHz}$ (\lr{Hz}) & \rowstrut\hspace{1.6cm} & \hspace{1.6cm} & \hspace{1.4cm} \\ \hline
		دومین انحرافی که حامل صفر می‌شود (\lr{Hz}) & \rowstrut & & \\ \hline
		پهنای باند با $\Delta f = 5\un{kHz}$ و $f_m = 1\un{kHz}$ (\lr{kHz}) & \rowstrut & & \\ \hline
		پهنای باند با $\Delta f = 5\un{kHz}$ و $f_m = 3\un{kHz}$ (\lr{kHz}) & \rowstrut & & \\ \hline
		$\mathrm{CNR}$ با \pv{tx\_atten}{40} (\lr{dB}) & \rowstrut & & \\ \hline
% 		خطای فرکانس میان فرستنده و گیرنده (\lr{Hz} و \lr{ppm}) & \rowstrut & & \\ \hline
	\end{tabularx}
\end{table}

\begin{expectbox}
حامل نخست در انحراف حدود
\q{2405}{Hz}
و سپس حدود
\q{5520}{Hz}
ناپدید می‌شود. خطای فرکانس میان دو دستگاه در این فرکانس می‌تواند چندین کیلوهرتز باشد؛ بدون تصحیح آن، سیگنال از فیلتر کانال
\q{16}{kHz}
بیرون می‌افتد.
\end{expectbox}

\begin{recordbox}
مقدار
\code{f_fine}،
$P_{\mathrm{off}}$،
$P_{\mathrm{on}}$،
تصویر طیف در لحظهٔ صفرشدن حامل و دو فایل
\IQ{}
برای پهنای باند.
\end{recordbox}

\begin{question}
\textbf{(اختیاری)}
انحراف‌های صفرشدن حامل را با مقدار نظری مقایسه کنید. اگر اختلافی هست، خطا از تنظیم
\code{dev}
است، از تفکیک نمایشگر یا از تشخیص «صفر»؟ برای هر کدام یک برآورد عددی بدهید.
\end{question}
\begin{question}
\textbf{(اختیاری)}
پهنای باند اندازه‌گیری‌شده را برای دو حالت با قاعدهٔ کارسون مقایسه کنید. در کدام حالت قاعدهٔ کارسون پهنای باند را کمتر از واقع برآورد می‌کند؟
\end{question}

% \labpart[۲۰ دقیقه]{نسبت سیگنال به نویز خروجی و آستانهٔ \lr{FM}}\label{part:thr}
% \begin{steps}
% \item فرستنده:
% \pv{dev}{5e3}،
% \pv{f\_m}{1e3}.
% گیرنده:
% \pv{dev}{5e3}،
% پهنای صوتی
% \q{3}{kHz}،
% \pv{pll\_bw}{0.05}.
% \item \code{tx_atten}
% را از
% \val{40}
% با گام
% \q{5}{dB}
% زیاد کنید تا
% \lr{SINAD}
% شروع به افت کند؛ از آن‌جا با گام
% \q{1}{dB}
% ادامه دهید تا
% \lr{SINAD}
% به کمتر از
% \q{5}{dB}
% برسد. اگر به انتهای بازهٔ
% \code{tx_atten}
% رسیدید،
% \code{tx_amp}
% را
% \val{0.05}
% کنید
% (\q{20}{dB}
% کمتر) و ادامه دهید.
% \item در هر پله: «توان کانال»،
% \lr{SINAD}
% برای هر دو آشکارساز و شنیده‌شدن یا نشدن «کلیک» را در جدول
% \ref{tab:thr}
% بنویسید و یک ضبط پنج‌ثانیه‌ای بگیرید. دست‌کم ده پله ثبت کنید که چهار تای آن در ناحیهٔ افت سریع باشد.
% \item $\mathrm{CNR}$
% هر پله را از
% $P_{\mathrm{on}}$
% و
% $P_{\mathrm{off}}$
% حساب کنید.
% \end{steps}

% \begin{table}[H]
% 	\centering\small
% 	\caption{$\mathrm{SNR}$ خروجی بر حسب $\mathrm{CNR}$ ($\Delta f = 5\un{kHz}$، $P_{\mathrm{off}} = $ \blank[1.3cm] \dBFS)}
% 	\label{tab:thr}
% 	\begin{tabularx}{\linewidth}{|c|Y|Y|Y|Y|Y|Y|}
% 		\hline
% 		\hdrrow \hdr{\lr{tx\_atten}} & \hdr{$P_{\mathrm{on}}$ (\dBFS)} & \hdr{$\mathrm{CNR}$ (\lr{dB})} & \hdr{\lr{SINAD} تمایزگر} & \hdr{\lr{SINAD} \lr{PLL}} & \hdr{کلیک؟} & \hdr{نام فایل} \\ \hline
% 		\rowstrut & & & & & & \\ \hline
% 		\rowstrut & & & & & & \\ \hline
% 		\rowstrut & & & & & & \\ \hline
% 		\rowstrut & & & & & & \\ \hline
% 		\rowstrut & & & & & & \\ \hline
% 		\rowstrut & & & & & & \\ \hline
% 		\rowstrut & & & & & & \\ \hline
% 		\rowstrut & & & & & & \\ \hline
% 		\rowstrut & & & & & & \\ \hline
% 		\rowstrut & & & & & & \\ \hline
% 	\end{tabularx}
% \end{table}

% \begin{expectbox}
% بالاتر از آستانه،
% \lr{SINAD}
% دسی‌بل به دسی‌بل با
% $\mathrm{CNR}$
% پایین می‌آید و حدود
% \q{13}{dB}
% از آن بیشتر است. در شبیه‌سازی با نویز ایده‌آل، برای
% $\Delta f = 5\un{kHz}$
% آستانه حوالی
% $\mathrm{CNR} \approx 8\un{dB}$
% است و دو آشکارساز تقریباً مانند هم رفتار می‌کنند؛ تفاوت آن‌ها با بزرگ‌شدن انحراف آشکار می‌شود (بخش بعد). در
% \lr{SINAD}های
% بالا عدد به سقفی می‌رسد که اعوجاج و نویز خود دستگاه‌ها تعیین می‌کند.
% \end{expectbox}

% \begin{question}
% در ناحیهٔ خطی، اختلاف
% \lr{SINAD}
% و
% $\mathrm{CNR}$
% را با رابطهٔ
% \eqref{eq:snr2}
% مقایسه کنید. آستانه را نقطه‌ای بگیرید که
% \lr{SINAD}
% یک دسی‌بل از خط راست پایین‌تر می‌افتد؛
% $\mathrm{CNR}$
% آن چقدر است؟
% \end{question}
% \begin{question}
% چرا برای محاسبهٔ
% $\mathrm{CNR}$
% باید
% $P_{\mathrm{off}}$
% را از
% $P_{\mathrm{on}}$
% کم کرد؟ این تصحیح در کدام ردیف‌های جدول مهم است؟
% \end{question}

% \labpart[۳۰ دقیقه]{چالش تیمی: انحراف فرکانس اختصاصی}\label{part:dev}
% \begin{tcolorbox}[sdrframed=sdrTeal,title={پارامتر اختصاصی تیم (دستیار آموزشی اعلام می‌کند)}]
% \small انحراف فرکانس تیم: $\Delta f_{\mathrm{team}} = $ \blank[2.2cm] \lr{kHz} \hfill تن آزمون: $f_m = 1\un{kHz}$ \hfill پهنای صوتی: $W = 3\un{kHz}$
% \end{tcolorbox}
% هدف: نشان دهید انحراف فرکانس تیم شما در مقایسه با
% $\Delta f = 5\un{kHz}$
% چه اثری بر پهنای باند، بر
% $\mathrm{SNR}$
% خروجی و بر آستانه دارد، و مزیت آشکارساز
% \lr{PLL}
% در آن چقدر است.
% \begin{steps}
% \item \textbf{پیش‌بینی.} پیش از اندازه‌گیری، ردیف «پیش‌بینی» جدول
% \ref{tab:devteam}
% را پر کنید: پهنای باند کارسون، تغییر
% \lr{SINAD}
% نسبت به
% \q{5}{kHz}
% در توان دریافتی یکسان، و مقدار مناسب
% \code{pll_bw}
% (با رابطهٔ
% \eqref{eq:recap2}
% و این قاعدهٔ سرانگشتی که
% $B_L$
% باید از هر دو مقدار
% $3W$
% و
% $4.7\sqrt{\Delta f\, f_m}$
% بزرگ‌تر باشد).
% \item \code{dev}
% را در فرستنده و گیرنده روی
% $\Delta f_{\mathrm{team}}$
% بگذارید؛ فیلتر کانال خودکار پهن یا باریک می‌شود.
% $P_{\mathrm{off}}$
% را دوباره بخوانید، چون پهنای کانال عوض شده است.
% \item پهنای باند را مانند بخش
% \ref{part:link}
% اندازه بگیرید و یک ضبط پنج‌ثانیه‌ای بگیرید.
% \item در سه سطح قوی که در بخش
% \ref{part:thr}
% هم اندازه گرفته‌اید،
% \lr{SINAD}
% را بخوانید و با مقادیر
% \q{5}{kHz}
% مقایسه کنید.
% \item ناحیهٔ آستانه را برای هر دو آشکارساز با گام
% \q{1}{dB}
% بپیمایید (دست‌کم شش پله، هر پله یک ضبط پنج‌ثانیه‌ای) و
% \code{tx_atten}
% آستانهٔ هر آشکارساز را بیابید.
% \item \code{pll_bw}
% را نصف و دو برابر کنید و اثر آن را بر آستانه و بر
% \lr{SINAD}
% در سیگنال قوی بنویسید.
% \end{steps}

% \begin{table}[H]
% 	\centering\small
% 	\caption{نتایج چالش انحراف فرکانس ($\Delta f_{\mathrm{team}} = $ \blank[1.4cm] \lr{kHz})}
% 	\label{tab:devteam}
% 	\begin{tabularx}{\linewidth}{|X|c|c|c|}
% 		\hline
% 		\hdrrow \hdr{کمیت} & \hdr{مرجع $5\un{kHz}$} & \hdr{پیش‌بینی برای تیم} & \hdr{اندازه‌گیری تیم} \\ \hline
% 		پهنای باند اشغالی (\lr{kHz}) & \rowstrut\hspace{1.8cm} & \hspace{1.8cm} & \hspace{1.8cm} \\ \hline
% 		\lr{SINAD} در سیگنال قوی، تمایزگر (\lr{dB}) & \rowstrut & & \\ \hline
% 		\code{tx_atten} آستانه، تمایزگر (\lr{dB}) & \rowstrut & & \\ \hline
% 		\code{tx_atten} آستانه، \lr{PLL} (\lr{dB}) & \rowstrut & & \\ \hline
% 		\code{pll_bw} به‌کاررفته و $B_L$ (\lr{kHz}) & \rowstrut & & \\ \hline
% 	\end{tabularx}
% \end{table}

% \begin{question}
% تغییر
% \lr{SINAD}
% نسبت به
% \q{5}{kHz}
% را با
% $20\log_{10}(\Delta f_{\mathrm{team}}/5\un{kHz})$
% مقایسه کنید. آستانهٔ تمایزگر (بر حسب توان دریافتی) چند دسی‌بل جابه‌جا شد و چرا؟
% \end{question}
% \begin{question}
% آشکارساز
% \lr{PLL}
% آستانه را چند دسی‌بل پایین‌تر برد؟ نتیجهٔ خود را با تیم‌هایی که انحراف کوچک‌تر و بزرگ‌تر داشتند مقایسه کنید و روند را توضیح دهید.
% \end{question}

%%%% SESSION-2 %%%%

%======================================================================
% \isection{flow}{فلوگراف‌های جلسهٔ دوم}
%======================================================================
% \subsection{گیرندهٔ استریو و \lr{RDS}: \lr{lab02\_stereo\_rds\_rx}}
% بخش ورودی این گیرنده همان گیرندهٔ جلسهٔ اول است (منبع، فیلتر انتخاب کانال با کاهش نرخ ۱۰ و
% \blk{Quadrature Demod}
% با
% $\Delta f = 75\un{kHz}$).
% خروجی تمایزگر، یعنی سیگنال
% \lr{MPX}
% با نرخ
% \q{240}{kHz}،
% به سه شاخه می‌رود (شکل
% \ref{fig:fg-stereo}):
% شاخهٔ مجموع، شاخهٔ پایلوت و تفاضل، و شاخهٔ
% \lr{RDS}.
% چون آشکارساز
% \lr{PLL}
% مؤلفه‌های بالای باند پایه را تضعیف می‌کند، در این گیرنده فقط تمایزگر به کار می‌رود.

% \begin{figure}[H]
% 	\centering
% 	\begin{tikzpicture}
% 		\node[sdrblkhl] (mpx) at (0,0) {MPX\\ 240 kHz};
% 		\node[sdrblk] (bpf) at (3.0,1.75) {\fa{میان‌گذر \lr{19 kHz}}\\ \fa{(ضرایب مختلط)}};
% 		\node[sdrblk] (pll) at (6.4,1.75) {PLL Carrier\\ Regeneration};
% 		\node[sdrblk] (sq) at (9.6,1.75) {\fa{$(\cdot)^2$ و بخش موهومی}};
% 		\node[sdrblk] (dl) at (3.0,0) {Delay};
% 		\node[sdrblk,circle,minimum width=0.75cm,minimum height=0.75cm,inner sep=0pt] (mx) at (6.4,0) {$\times$};
% 		\node[sdrblk] (dif) at (9.6,0) {\fa{پایین‌گذر \lr{15 kHz}، $\downarrow 5$}\\ \fa{بهره: \lr{lr\_gain}}};
% 		\node[sdrblk] (sum) at (9.6,-1.5) {\fa{پایین‌گذر \lr{15 kHz}، $\downarrow 5$}};
% 		\node[sdrblkpc] (mat) at (13.2,-0.75) {$L = \Sigma + \Delta$\\ $R = \Sigma - \Delta$\\ \fa{واتأکید و صدا}};
% 		\draw[flt] (mpx.north) |- (bpf.west); \draw[flt] (mpx)--(dl); \draw[cplx] (bpf)--(pll); \draw[cplx] (pll)--(sq);
% 		\draw[flt] (dl)--(mx); \draw[flt] (sq.south) -- ++(0,-0.45) -| node[pos=0.25,above,font=\scriptsize,text=black]{38 kHz}(mx.north);
% 		\draw[flt] (mx)--(dif); \draw[flt] (dl.south) |- (sum.west);
% 		\draw[flt] (dif.east) -- ++(0.35,0) |- ([yshift=4pt]mat.west); \draw[flt] (sum.east) -- ++(0.35,0) |- ([yshift=-4pt]mat.west);
% 		\node[font=\scriptsize,text=sdrGray] at (11.45,0.25) {$L-R$}; \node[font=\scriptsize,text=sdrGray] at (11.45,-1.75) {$L+R$};
% 		\node[sdrblk] (x57) at (2.2,-3.3) {\fa{جابه‌جایی \lr{57 kHz}}\\ $\downarrow 10$};
% 		\node[sdrblk] (rs) at (5.0,-3.3) {\fa{تغییر نرخ}\\ $\times 19/24$};
% 		\node[sdrblk] (mf) at (7.75,-3.3) {\fa{فیلتر منطبق}\\ \fa{و \lr{AGC}}};
% 		\node[sdrblk] (sy) at (10.6,-3.3) {Symbol Sync\\ \fa{و گیرندهٔ \lr{BPSK}}};
% 		\node[sdrblkpc] (rd) at (13.5,-3.3) {RDS Decoder\\ Parser, Panel};
% 		\draw[flt] (mpx.south) |- (x57.west);
% 		\draw[cplx] (x57)--node[above,font=\scriptsize,text=black]{24k}(rs); \draw[cplx] (rs)--node[above,font=\scriptsize,text=black]{19k}(mf);
% 		\draw[cplx] (mf)--(sy); \draw[msg] (sy)--node[above,font=\scriptsize,text=black]{\fa{بیت}}(rd);
% 	\end{tikzpicture}
% 	\caption{ساختار گیرندهٔ استریو و \lr{RDS}: بالا دیکودر استریو، پایین زنجیرهٔ \lr{RDS} (رسم‌شده؛ تصویر واقعی \GRC{} نیست)}
% 	\label{fig:fg-stereo}
% \end{figure}
% \screenshotslot{labs/figs/lab02_stereo_rds_rx.png}{تصویر فلوگراف \lr{lab02\_stereo\_rds\_rx} در \GRC}

% \begin{table}[H]
% 	\centering\small
% 	\caption{پارامترهای دیکودر استریو (ساختار برگرفته از نمونهٔ \lr{gr-rds}؛ نام‌ها مطابق \GR{} \lr{3.10})}
% 	\label{tab:params-st}
% 	\begin{tabularx}{\linewidth}{|C{4.4cm}|Z|}
% 		\hline
% 		\hdrrow \hdr{بلوک (\lr{Id})} & \hdr{پارامترها} \\ \hline
% 		\blk{Quadrature Demod} & \pv{Gain}{quad\_rate/(2*math.pi*75e3)} \\ \hline
% 		\blk{Band-pass Filter Taps} (\code{pilot_taps}) & \pv{Tap Type}{Complex}؛ \pv{Sample Rate (Hz)}{quad\_rate}؛ \pv{Low Cutoff Freq (Hz)}{18980}؛ \pv{High Cutoff Freq (Hz)}{19020}؛ \pv{Transition Width (Hz)}{1000} \\ \hline
% 		\blk{Decimating FIR Filter} (پایلوت) & \pv{Type}{Float->Complex (Complex Taps)}؛ \pv{Decimation}{1}؛ \pv{Taps}{pilot\_taps} \\ \hline
% 		\blk{PLL Carrier}\newline\blk{Regeneration} & \pv{Loop Bandwidth}{0.001}؛ \pv{Max Phase/sample}{2*math.pi*19020/quad\_rate}؛ \pv{Min Phase/sample}{2*math.pi*18980/quad\_rate} \\ \hline
% 		\blk{Multiply} و \blk{Complex to Imag} & \blk{Multiply} مختلط با دو ورودی که هر دو به خروجی \lr{PLL} وصل‌اند (مربع پایلوت)، سپس بخش موهومی: حامل \q{38}{kHz} \\ \hline
% 		\blk{Delay} & \pv{Type}{float}؛ \pv{Delay}{(len(pilot\_taps) - 1) // 2}؛ ورودی: خروجی تمایزگر \\ \hline
% 		\blk{Decimating FIR Filter} (مجموع) & \pv{Type}{Float->Float (Real Taps)}؛ \pv{Decimation}{5}؛ \pv{Taps}{firdes.low\_pass(1.0, quad\_rate, 15e3, 2e3)}؛ ورودی: خروجی \blk{Delay} \\ \hline
% 		\blk{Multiply} (حقیقی) و \blk{Decimating FIR Filter} (تفاضل) & ضرب خروجی \blk{Delay} در حامل \q{38}{kHz}؛ سپس \pv{Decimation}{5}؛ \pv{Taps}{firdes.low\_pass(lr\_gain*stereo, quad\_rate, 15e3, 2e3)} \\ \hline
% 		\blk{Add} و \blk{Subtract} & \blk{Add}: مجموع و تفاضل (کانال چپ)؛ \blk{Subtract}: مجموع منهای تفاضل (کانال راست) \\ \hline
% 		\blk{FM Deemphasis}، \blk{Multiply Const}، \blk{Audio Sink} & برای هر کانال \pv{Tau}{tau} و ضرب در \code{volume}؛ \blk{Audio Sink} با \pv{Num Inputs}{2} \\ \hline
% 		\blk{QT GUI Range}، \blk{QT GUI Chooser} & \code{lr_gain} (\val{-2.1}، از \val{-3} تا \val{-1}، گام \val{0.05})؛ \code{stereo}: \val{1} (استریو) یا \val{0} (مونو) \\ \hline
% 	\end{tabularx}
% \end{table}
% \begin{table}[H]
% 	\centering\small
% 	\caption{پارامترهای زنجیرهٔ \lr{RDS} و نمایشگرها}
% 	\label{tab:params-rds}
% 	\begin{tabularx}{\linewidth}{|C{4.4cm}|Z|}
% 		\hline
% 		\hdrrow \hdr{بلوک (\lr{Id})} & \hdr{پارامترها} \\ \hline
% 		\blk{Frequency Xlating}\newline\blk{FIR Filter} & \pv{Type}{Float->Complex (Complex Taps)}؛ \pv{Decimation}{10}؛ \pv{Taps}{firdes.low\_pass(1.0, quad\_rate, 7.5e3, 5e3)}؛ \pv{Center Frequency}{57e3}؛ \pv{Sample Rate}{quad\_rate} \\ \hline
% 		\blk{Rational Resampler} & \pv{Type}{Complex->Complex (Complex Taps)}؛ \pv{Interpolation}{19}؛ \pv{Decimation}{24} \\ \hline
% 		\blk{RRC Filter Taps} (\code{rrc_taps}) و \blk{Variable} & \pv{Sample Rate (Hz)}{19000}؛ \pv{Symbol Rate (Hz)}{19000/8}؛ \pv{Excess BW}{1.0}؛ \pv{Num Taps}{151}. سپس متغیر \code{rrc_taps_manchester} با این مقدار (یک عبارت پیوسته): \val{[rrc\_taps[n] - rrc\_taps[n+8]} \val{for n in range(len(rrc\_taps)-8)]} \\ \hline
% 		\blk{Decimating FIR Filter} (منطبق) & \pv{Type}{Complex->Complex (Complex Taps)}؛ \pv{Decimation}{1}؛ \pv{Taps}{rrc\_taps\_manchester} \\ \hline
% 		\blk{AGC} & \pv{Type}{complex}؛ \pv{Rate}{2e-3}؛ \pv{Reference}{0.585}؛ \pv{Gain}{53}؛ \pv{Max Gain}{1000} \\ \hline
% 		\blk{Symbol Sync} & \pv{I/O Type}{Complex}؛ \pv{Timing Error Detector}{Zero Crossing}؛ \pv{TED Slicer Constellation}{digital.constellation\_bpsk().base()}؛ \pv{Samples per Symbol}{16}؛ \pv{Loop Bandwidth}{0.01}؛ \pv{Maximum Deviation}{0.1}؛ \pv{Output Samples/Symbol}{1} \\ \hline
% 		\blk{Constellation Receiver} & \pv{Constellation Object}{digital.constellation\_bpsk().base()}؛ \pv{Loop Bandwidth}{2*math.pi/100}؛ \pv{Minimum Freq Deviation}{-0.002}؛ \pv{Maximum Freq Deviation}{0.002}. خروجی صفر به \blk{Differential Decoder}، خروجی‌های ۱ تا ۳ به \blk{Null Sink} و خروجی ۴ به \blk{QT GUI Constellation Sink} \\ \hline
% 		\blk{Differential Decoder} & \pv{Modulus}{2} \\ \hline
% 		\blk{RDS Decoder}، \blk{RDS Parser}، \blk{RDS Panel} & اتصال پیامی (درگاه‌های خاکستری): \lr{out} دیکودر به \lr{in} پارسر و \lr{out} پارسر به \lr{in} پنل \\ \hline
% 		\blk{Python Block} (شمارندهٔ گروه) & کد \code{sdrlab_rds_counter.py} (پیوست)؛ \pv{Period}{5.0}؛ درگاه \lr{in} به \lr{out} بلوک \blk{RDS Decoder} \\ \hline
% 		نمایشگرها & \blk{QT GUI Frequency Sink} حقیقی روی \lr{MPX} (\pv{FFT Size}{4096}؛ \pv{Bandwidth (Hz)}{quad\_rate})؛ \blk{QT GUI Frequency Sink} حقیقی با دو ورودی روی کانال چپ و راست (\pv{Bandwidth (Hz)}{audio\_rate}) \\ \hline
% 	\end{tabularx}
% \end{table}

% \begin{notebox}[ضریب \lr{lr\_gain}]
% علامت منفی، فاز حامل بازسازی‌شده را تصحیح می‌کند؛ عدد ۲ نصف‌شدن دامنه در آشکارسازی
% \lr{DSB-SC}
% را جبران می‌کند و مقدار اضافه (از ۲ تا حدود
% $2.1$)
% افت دامنهٔ زیرحامل در فیلتر کانال را. بهترین مقدار را در بخش
% \ref{part:stereo}
% با اندازه‌گیری می‌یابید.
% \end{notebox}

% \subsection{فرستندهٔ آمادهٔ استریو و \lr{RDS}: \lr{lab02\_stereo\_rds\_tx}}
% این فلوگراف را آماده دریافت می‌کنید و فقط متغیرهای جدول
% \ref{tab:params-sttx}
% را تغییر می‌دهید. سیگنال
% \lr{MPX}
% در آن با نرخ
% \q{380}{kHz}
% ساخته می‌شود (۲۰ برابر پایلوت، تا هر نیم‌بیت
% \lr{RDS}
% تعداد صحیحی نمونه داشته باشد) و پس از مدولاسیون با نسبت
% $120/19$
% به
% \q{2.4}{MS/s}
% می‌رسد. پایلوت، حامل
% \q{38}{kHz}
% و حامل
% \q{57}{kHz}
% همه از یک نوسان‌ساز
% \q{19}{kHz}
% ساخته می‌شوند و بنابراین هم‌فازند. فلوگراف را باز کنید و مسیر هر جزء رابطهٔ
% \lr{MPX}
% را در آن پیدا کنید.

% \begin{table}[H]
% 	\centering\small
% 	\caption{متغیرهای قابل تنظیم فرستندهٔ آماده}
% 	\label{tab:params-sttx}
% 	\begin{tabularx}{\linewidth}{|c|c|X|}
% 		\hline
% 		\hdrrow \hdr{متغیر} & \hdr{مقدار آغازین} & \hdr{معنا} \\ \hline
% 		\code{f_tx}، \code{tx_atten}، \code{tx_amp} & \val{434.0e6}، \val{40}، \val{0.5} & مانند فرستندهٔ جلسهٔ اول \\ \hline
% 		\code{f_L}، \code{a_L} & \val{1e3}، \val{0.9} & فرکانس و دامنهٔ تن کانال چپ \\ \hline
% 		\code{f_R}، \code{a_R} & \val{2.5e3}، \val{0.0} & فرکانس و دامنهٔ تن کانال راست \\ \hline
% 		\code{pilot_gain} & \val{0.1} & سطح پایلوت نسبت به بیشینهٔ انحراف \\ \hline
% 		\code{rds_gain} & \val{0.027} & سطح \lr{RDS}؛ \val{0.027} یعنی انحراف حدود \q{2}{kHz} \\ \hline
% 		\code{ps_name}، \code{pi_ref} & \val{"TEAM-00"}، \val{0} & نام ایستگاه (حداکثر هشت حرف) و شمارهٔ مرجع برنامه؛ شناسهٔ برنامه \lr{D3} به‌اضافهٔ \code{pi_ref} به‌صورت دو رقم شانزده‌شانزدهی می‌شود \\ \hline
% 	\end{tabularx}
% \end{table}

%======================================================================
% \isection{steps}{شرح آزمایش: جلسهٔ دوم}
%======================================================================

% \labpart[۲۵ دقیقه]{گیرندهٔ \lr{WFM} مونو و طیف باند پایه}\label{part:mono}
% \begin{steps}
% \item آنتن را ببندید (چیدمان الف) و فلوگراف
% \code{lab02_stereo_rds_rx.grc}
% را در حالت
% \lr{Mono}
% روی یک ایستگاه قوی اجرا کنید. بهره را بدون اشباع تنظیم کنید.
% \item در نمایشگر
% \lr{MPX}
% با
% \pv{FFT Size}{4096}
% (فاصلهٔ خانه‌ها
% \q{58.6}{Hz})
% و میانگین‌گیری
% \val{High}،
% اجزای طیف باند پایه را بیابید و جدول
% \ref{tab:mpx}
% را پر کنید. برای خواندن دقیق فرکانس، روی هر جزء بزرگ‌نمایی کنید.
% \item همین کار را برای یک ایستگاه دیگر تکرار کنید و تفاوت‌ها را بنویسید (آیا پایلوت دارد؟ آیا
% \lr{RDS}
% دارد؟).
% \item ثابت زمانی واتأکید را میان
% \q{50}{}
% و
% \q{75}{$\mu$s}
% عوض کنید و اثر آن را بر صدا و بر شیب طیف صوت ببینید.
% \item یک ضبط ده‌ثانیه‌ای از ایستگاه نخست بگیرید.
% \end{steps}

% \begin{table}[H]
% 	\centering\small
% 	\caption{اجزای طیف باند پایهٔ استریو (ایستگاه: \blank[1.8cm] \lr{MHz}، بهره: \blank[1.2cm] \lr{dB})}
% 	\label{tab:mpx}
% 	\begin{tabularx}{\linewidth}{|X|c|c|c|}
% 		\hline
% 		\hdrrow \hdr{جزء} & \hdr{فرکانس مورد انتظار (\lr{kHz})} & \hdr{فرکانس اندازه‌گیری‌شده (\lr{kHz})} & \hdr{سطح (\lr{dB})} \\ \hline
% 		لبهٔ بالایی $L+R$ & $15$ & \rowstrut\hspace{2.2cm} & \hspace{1.5cm} \\ \hline
% 		پایلوت & $19.000$ & \rowstrut & \\ \hline
% 		لبهٔ پایینی و بالایی $L-R$ & $23$ و $53$ & \rowstrut & \\ \hline
% 		مرکز و پهنای \lr{RDS} & $57$ و حدود $4.8$ & \rowstrut & \\ \hline
% 		کف نویز حوالی \q{17}{kHz} & — & — & \rowstrut \\ \hline
% 		کف نویز حوالی \q{65}{kHz} & — & — & \rowstrut \\ \hline
% 	\end{tabularx}
% \end{table}

% \begin{expectbox}
% پایلوت خطی باریک در
% \q{19}{kHz}
% است؛ دو نیمهٔ
% $L-R$
% دو طرف
% \q{38}{kHz}
% قرینهٔ هم‌اند و در خود
% \q{38}{kHz}
% حاملی دیده نمی‌شود. ایستگاهی که
% \lr{RDS}
% دارد برآمدگی کوچکی به پهنای حدود
% \q{5}{kHz}
% در
% \q{57}{kHz}
% نشان می‌دهد. کف نویز با فرکانس بالا می‌رود. همهٔ ایستگاه‌ها استریو یا دارای
% \lr{RDS}
% نیستند.
% \end{expectbox}
% \screenshotslot{labs/figs/lab02_mpx_spectrum.png}{نمونهٔ طیف باند پایهٔ یک ایستگاه استریو در آزمایشگاه شما}

% \begin{question}
% چرا کف نویز طیف باند پایه با فرکانس بالا می‌رود؟ اختلاف سطح نویز در
% \q{17}{}
% و
% \q{65}{kHz}
% را با پیش‌بینی
% $f^2$
% مقایسه کنید.
% \end{question}
% \begin{question}
% فرکانس پایلوت را با چه دقتی خواندید و این دقت را چه چیزی محدود می‌کند؟ خطای نوسان‌ساز گیرنده
% (چند ده
% \lr{ppm})
% چه سهمی در آن دارد؟
% \end{question}

% \labpart[۳۵ دقیقه]{بازسازی استریو}\label{part:stereo}
% \begin{steps}
% \item حالت را روی
% \lr{Stereo}
% بگذارید و یک ایستگاه استریو را با هدفون بشنوید. میان مونو و استریو چند بار جابه‌جا شوید؛ روی یک ایستگاه ضعیف‌تر، تفاوت نویز دو حالت را نیز بشنوید.
% \item \textbf{آزمون با فرستندهٔ کابلی.} فلوگراف‌ها را متوقف کنید، چیدمان (ب) را ببندید و چک‌لیست راه‌اندازی را کامل کنید. فلوگراف
% \code{lab02_stereo_rds_tx.grc}
% را با مقادیر آغازین جدول
% \ref{tab:params-sttx}
% اجرا کنید (فقط کانال چپ تن دارد). گیرنده را روی
% $f_{\mathrm{tx}}$
% بگذارید.
% \item در نمایشگر «کانال چپ و راست» سطح تن
% \q{1}{kHz}
% را در هر دو کانال بخوانید؛ اختلاف این دو، جداسازی استریو است. سپس
% \pv{a\_L}{0}
% و
% \pv{a\_R}{0.9}
% بگذارید و جداسازی را در جهت دیگر اندازه بگیرید.
% \item \code{lr_gain}
% را در بازهٔ
% \val{-1.8}
% تا
% \val{-2.4}
% تغییر دهید و مقداری را که بیشترین جداسازی را می‌دهد بیابید (جدول
% \ref{tab:sep}).
% \item \textbf{اثر تأخیر.} مقدار
% \blk{Delay}
% را یک، دو و سه نمونه بیشتر کنید (هر بار فلوگراف را دوباره اجرا کنید) و جداسازی را بخوانید. سپس مقدار درست را بازگردانید.
% \item \textbf{اثر پایلوت.} در فرستنده
% \code{pilot_gain}
% را صفر کنید. خروجی دو کانال چه می‌شود؟ سپس آن را به
% \val{0.1}
% بازگردانید.
% \item یک ضبط ده‌ثانیه‌ای از سیگنال کابلی بگیرید.
% \end{steps}

% \begin{table}[H]
% 	\centering\small
% 	\caption{جداسازی استریو با فرستندهٔ کابلی (تن فقط در کانال چپ مگر خلاف آن ذکر شود)}
% 	\label{tab:sep}
% 	\begin{tabularx}{\linewidth}{|X|Y|Y|Y|}
% 		\hline
% 		\hdrrow \hdr{تنظیم} & \hdr{سطح تن در چپ (\lr{dB})} & \hdr{سطح تن در راست (\lr{dB})} & \hdr{جداسازی (\lr{dB})} \\ \hline
% 		\pv{lr\_gain}{-2.1}، تأخیر درست & \rowstrut & & \\ \hline
% 		تن فقط در کانال راست & \rowstrut & & \\ \hline
% 		\pv{lr\_gain}{-2.0} & \rowstrut & & \\ \hline
% 		بهترین \code{lr_gain} = \blank[1.2cm] & \rowstrut & & \\ \hline
% 		تأخیر یک نمونه بیشتر & \rowstrut & & \\ \hline
% 		تأخیر دو نمونه بیشتر & \rowstrut & & \\ \hline
% 		تأخیر سه نمونه بیشتر & \rowstrut & & \\ \hline
% 		\pv{pilot\_gain}{0} & \rowstrut & & \\ \hline
% 	\end{tabularx}
% \end{table}

% \begin{expectbox}
% در شبیه‌سازی با کانال ایده‌آل، جداسازی با
% \pv{lr\_gain}{-2.1}
% حدود
% \q{49}{dB}
% و با
% \val{-2.0}
% حدود
% \q{34}{dB}
% است؛ روی سخت‌افزار معمولاً کمتر است. یک نمونه خطای تأخیر جداسازی را به حدود
% \q{11}{dB}
% می‌رساند و با سه نمونه خطا دو کانال تقریباً جابه‌جا می‌شوند، چون سه نمونه در نرخ
% \q{240}{kHz}
% نزدیک نیم‌دورهٔ حامل
% \q{38}{kHz}
% است.
% \end{expectbox}

% \begin{question}
% چرا بلوک
% \blk{Delay}
% لازم است و مقدار آن از کجا می‌آید؟ جداسازی اندازه‌گیری‌شده با یک نمونه خطا را با رابطهٔ خطای فاز در گزارش پیش‌از‌جلسه مقایسه کنید.
% \end{question}
% \begin{question}
% با حذف پایلوت چه رخ داد و چرا؟ گیرنده‌های تجاری در این وضعیت چه می‌کنند؟
% \end{question}
% \begin{question}
% روی ایستگاه ضعیف، نویز حالت استریو از مونو بیشتر بود. با توجه به شکل طیف نویز خروجی تمایزگر توضیح دهید چرا.
% \end{question}

% \labpart[۲۵ دقیقه]{دیکد \lr{RDS}}\label{part:rds}
% \begin{steps}
% \item فرستندهٔ کابلی هنوز وصل است. در فرستنده
% \code{ps_name}
% را شناسهٔ تیم و
% \code{pi_ref}
% را شمارهٔ تیم بگذارید و دوباره اجرا کنید
% (\pv{rds\_gain}{0.027}).
% \item در گیرنده نمودار صورت‌فلکی
% \lr{RDS}
% را ببینید: پس از چند ثانیه دو خوشه روی محور حقیقی شکل می‌گیرد. در
% \blk{RDS Panel}
% نام ایستگاه و شناسهٔ برنامه را بخوانید و در کنسول، سطرهای
% \code{[RDS]}
% شمارندهٔ گروه را ببینید.
% \item شناسهٔ برنامهٔ نمایش‌داده‌شده را با مقدار مورد انتظار
% (\lr{D3}
% و شمارهٔ تیم) مقایسه کنید.
% \item فلوگراف فرستنده را متوقف کنید، کابل را جدا و آنتن را ببندید. دست‌کم چهار ایستگاه را بررسی و جدول
% \ref{tab:stations}
% را پر کنید.
% \item برای یک ایستگاه دارای
% \lr{RDS}،
% روی
% \q{57}{kHz}
% در طیف
% \lr{MPX}
% بزرگ‌نمایی و پهنای برآمدگی
% \lr{RDS}
% و سطح آن نسبت به پایلوت را اندازه بگیرید. یک ضبط بیست‌ثانیه‌ای بگیرید.
% \end{steps}

% \begin{table}[H]
% 	\centering\small
% 	\caption{ایستگاه‌های بررسی‌شده}
% 	\label{tab:stations}
% 	\begin{tabularx}{\linewidth}{|c|c|c|c|Y|Y|Y|}
% 		\hline
% 		\hdrrow \hdr{فرکانس (\lr{MHz})} & \hdr{توان کانال (\dBFS)} & \hdr{پایلوت؟} & \hdr{\lr{RDS}؟} & \hdr{نام ایستگاه (\lr{PS})} & \hdr{شناسه (\lr{PI})} & \hdr{گروه در \q{5}{s}} \\ \hline
% 		\rowstrut\hspace{1.6cm} & \hspace{1.6cm} & & & & & \\ \hline
% 		\rowstrut & & & & & & \\ \hline
% 		\rowstrut & & & & & & \\ \hline
% 		\rowstrut & & & & & & \\ \hline
% 		\rowstrut & & & & & & \\ \hline
% 	\end{tabularx}
% \end{table}

% \begin{expectbox}
% با سیگنال قوی، شمارنده حدود ۵۷ گروه در هر پنج ثانیه نشان می‌دهد (نرخ اسمی
% $11.4$
% گروه بر ثانیه) و نام ایستگاه ظرف یک تا دو ثانیه پس از قفل کامل می‌شود. اگر هیچ ایستگاهی در محل شما
% \lr{RDS}
% ندارد، این بخش را فقط با فرستندهٔ کابلی انجام دهید و همین را در گزارش بنویسید.
% \end{expectbox}

% \begin{question}
% چرا شاخهٔ
% \lr{RDS}
% به نرخ
% \q{19}{kHz}
% برده می‌شود؟ در این نرخ هر بیت و هر نیم‌بیت چند نمونه دارد و چرا این برای بلوک
% \blk{Symbol Sync}
% مهم است؟
% \end{question}
% \begin{question}
% از لحظهٔ اجرای فلوگراف تا نمایش کامل نام ایستگاه چقدر طول کشید؟ این زمان را به سهم همگام‌شدن و سهم ساختار گروه‌ها تفکیک کنید.
% \end{question}

% \labpart[۳۰ دقیقه]{چالش تیمی: \lr{RDS} در برابر سطح سیگنال}\label{part:rdslevel}
% \begin{tcolorbox}[sdrframed=sdrTeal,title={پارامترهای اختصاصی تیم (دستیار آموزشی اعلام می‌کند)}]
% \small سطح \lr{RDS}: \code{rds_gain} = \blank[1.8cm] \hfill نام ایستگاه: \code{ps_name} = \blank[2.4cm] \hfill شمارهٔ مرجع: \code{pi_ref} = \blank[1.2cm]
% \end{tcolorbox}
% هدف: منحنی نرخ خطای گروه
% \lr{RDS}
% را بر حسب
% $\mathrm{CNR}$
% برای سطح
% \lr{RDS}
% تیم خود به دست آورید و با کیفیت صدا در همان سطح‌ها مقایسه کنید.
% \begin{steps}
% \item چیدمان (ب) را ببندید. فرستنده را با پارامترهای تیم و
% \pv{a\_L}{0.5}،
% \pv{a\_R}{0.5}
% اجرا کنید.
% \item فلوگراف فرستنده را متوقف و «توان کانال» را بخوانید
% ($P_{\mathrm{off}}$
% در پهنای
% \q{200}{kHz})؛
% سپس دوباره اجرا کنید.
% \item \code{tx_atten}
% را از
% \val{40}
% با گام
% \q{5}{dB}
% زیاد کنید تا شمارندهٔ گروه شروع به کم شدن کند؛ از آن‌جا با گام
% \q{1}{dB}
% ادامه دهید تا هیچ گروهی دیکد نشود. در صورت نیاز
% \code{tx_amp}
% را ده برابر کوچک کنید.
% \item در هر پله سه سطر پیاپی شمارنده را
% (\q{15}{s})
% بخوانید و میانگین بگیرید. نرخ خطای گروه برابر
% $1 - n/57$
% است که
% $n$
% میانگین تعداد گروه در پنج ثانیه است. کیفیت صدای مونو و استریو را نیز در همان پله ارزیابی کنید (جدول
% \ref{tab:rdslevel}).
% \item در سه پله (بالا، میانه و پایین ناحیهٔ گذار) یک ضبط بیست‌ثانیه‌ای بگیرید.
% \item $\mathrm{CNR}$
% متناظر با نرخ خطای ۵۰ درصد را با درون‌یابی به دست آورید.
% \end{steps}

% \begin{table}[H]
% 	\centering\small
% 	\caption{نرخ خطای گروه \lr{RDS} بر حسب سطح سیگنال (\lr{rds\_gain} = \blank[1.2cm] ، $P_{\mathrm{off}} = $ \blank[1.3cm] \dBFS)}
% 	\label{tab:rdslevel}
% 	\begin{tabularx}{\linewidth}{|c|Y|Y|Y|Y|Y|Y|}
% 		\hline
% 		\hdrrow \hdr{\lr{tx\_atten}} & \hdr{$P_{\mathrm{on}}$ (\dBFS)} & \hdr{$\mathrm{CNR}$ (\lr{dB})} & \hdr{گروه در \q{5}{s}} & \hdr{نرخ خطای گروه (٪)} & \hdr{نام درست؟} & \hdr{صدا: مونو / استریو} \\ \hline
% 		\rowstrut & & & & & & \\ \hline
% 		\rowstrut & & & & & & \\ \hline
% 		\rowstrut & & & & & & \\ \hline
% 		\rowstrut & & & & & & \\ \hline
% 		\rowstrut & & & & & & \\ \hline
% 		\rowstrut & & & & & & \\ \hline
% 		\rowstrut & & & & & & \\ \hline
% 		\rowstrut & & & & & & \\ \hline
% 		\rowstrut & & & & & & \\ \hline
% 	\end{tabularx}
% \end{table}

% \begin{expectbox}
% در شبیه‌سازی با نویز ایده‌آل و
% \pv{rds\_gain}{0.027}،
% دیکد در
% $\mathrm{CNR}$
% حدود
% \q{16}{dB}
% تقریباً از کار می‌افتد و از حدود
% \q{21}{dB}
% به بالا تقریباً بی‌خطاست؛ دو برابر کردن
% \code{rds_gain}
% این ناحیه را حدود
% \q{6}{dB}
% پایین می‌برد. گذار بسیار تند است: چند دسی‌بل میان «همه» و «هیچ».
% \end{expectbox}

% \begin{question}
% جابه‌جایی ناحیهٔ گذار تیم شما نسبت به
% \pv{rds\_gain}{0.027}
% را با
% $20\log_{10}(\text{\lr{rds\_gain}}/0.027)$
% مقایسه کنید. نتیجهٔ تیم‌های دیگر این رابطه را تأیید می‌کند؟
% \end{question}
% \begin{question}
% در سطحی که
% \lr{RDS}
% از کار افتاد، صدای مونو و استریو چه وضعی داشتند؟ با توجه به شکل طیف نویز و سهم انحراف هر جزء توضیح دهید چرا
% \lr{RDS}
% زودتر از صدا از بین می‌رود.
% \end{question}

%======================================================================
\isection{wrench}{خطایابی}
%======================================================================
\begin{table}[H]
	\centering\small
	\caption{نشانه‌ها، علت‌های محتمل و راه‌حل‌ها}
	\label{tab:trouble-fm}
	\begin{tabularx}{\linewidth}{|X|X|X|}
		\hline
		\hdrrow \hdr{نشانه} & \hdr{علت محتمل} & \hdr{راه‌حل} \\ \hline
		فقط خش‌خش شنیده می‌شود & فرکانس اشتباه، یا \code{dev} و پهنای صوتی با سیگنال نمی‌خواند & \code{fc} را بررسی کنید؛ برای پخش \lr{FM}: \pv{dev}{75e3} و پهنای صوتی \q{15}{kHz} \\ \hline
		صدا بریده‌بریده است؛ \lr{aU} در کنسول & کارت صدا و گیرنده دو ساعت جدا دارند؛ یا بار پردازنده زیاد است & گاه‌به‌گاه طبیعی است؛ برنامه‌های دیگر را ببندید؛ \pv{OK to Block}{Yes} \\ \hline
		صدا خفه یا بیش از حد تیز است & ثابت زمانی واتأکید نادرست & \code{tau} را عوض کنید \\ \hline
		سیگنال کابلی در کانال دیده نمی‌شود & خطای فرکانس از پهنای فیلتر کانال بیشتر است؛ یا فرستنده اجرا نمی‌شود & موقتاً \code{dev} گیرنده را بزرگ کنید، سیگنال را بیابید و با \code{f_fine} به مرکز بیاورید \\ \hline
		\lr{SINAD} حتی با سیگنال قوی کم است & \code{f_m} گیرنده با تن فرستنده یکی نیست؛ خطای فرکانس؛ یا اشباع گیرنده & \code{f_m} و \code{dev} دو طرف را یکی کنید؛ خطا را صفر کنید؛ اگر توان کانال نزدیک صفر \dBFS{} است تضعیف را زیاد کنید \\ \hline
		خروجی \lr{PLL} اعوجاج دارد & حلقه باریک‌تر از آن است که انحراف را دنبال کند & \code{pll_bw} را زیاد کنید \\ \hline
		فرستنده پیدا نمی‌شود & نشانی دستگاه یا درایور & در \blk{PlutoSDR Sink} مقدار \prm{IIO context URI} را بررسی کنید (مثلاً \val{ip:192.168.2.1})؛ از دستیار آموزشی کمک بگیرید \\ \hline
		قله‌ای باریک در مرکز طیف فرستنده & نشت حامل فرستنده & طبیعی است؛ به همین دلیل سیگنال با \code{f_off} از مرکز دور شده است \\ \hline
% 		دو کانال صدا یکسان‌اند & حالت \lr{Mono}؛ یا شاخهٔ پایلوت قفل نیست & حالت را \lr{Stereo} کنید؛ حدود فرکانس \lr{PLL} پایلوت و اتصال \blk{Multiply} را بررسی کنید \\ \hline
% 		کانال چپ و راست جابه‌جا شده‌اند & علامت \code{lr_gain} یا مقدار \blk{Delay} نادرست است & \code{lr_gain} منفی باشد؛ \blk{Delay} را طبق جدول \ref{tab:params-st} بگذارید \\ \hline
% 		\blk{RDS Panel} خالی می‌ماند & ایستگاه \lr{RDS} ندارد؛ سیگنال ضعیف است؛ نسبت تغییر نرخ نادرست؛ یا اتصال پیامی برقرار نیست & برآمدگی \q{57}{kHz} را در طیف \lr{MPX} بیابید؛ نسبت $19/24$ و اتصال درگاه‌های خاکستری را بررسی کنید \\ \hline
		خطای \lr{ID is blacklisted} در \GRC & شناسهٔ بلوک با یک نام رزروشده یکی است (مثلاً \code{audio}) & شناسهٔ بلوک را عوض کنید \\ \hline
	\end{tabularx}
\end{table}

%======================================================================
\isection{pencil}{تحلیل نتایج و گزارش پس‌از‌جلسه}
%======================================================================
گزارش در قالب اجباری آیین‌نامه (بند
\bnd{۴-۲})
و بر پایهٔ فایل‌های خام ذخیره‌شده نوشته می‌شود. مهلت و این‌که برای این آزمایش دوجلسه‌ای یک گزارش تحویل شود یا دو گزارش را مدرس اعلام می‌کند؛ تحلیل‌های الزامی هر جلسه در زیر جدا آمده است.

\subsection{تحلیل‌های الزامی جلسهٔ اول}
\begin{enumerate}
% \item \textbf{منحنی $\mathrm{SNR}$ خروجی بر حسب $\mathrm{SNR}$ ورودی.} از فایل‌های بخش
% \ref{part:thr}
% و با کد خودتان (یا با اجرای فلوگراف روی
% \blk{File Source})،
% $\mathrm{SNR}$
% خروجی صوت را برای هر دو آشکارساز بر حسب
% $\mathrm{CNR}$
% رسم و نقطهٔ آستانهٔ
% \lr{FM}
% را به‌طور تجربی تعیین کنید. منحنی را با رابطهٔ
% \eqref{eq:snr2}
% و با شبیه‌سازی پیش‌از‌جلسه مقایسه کنید.
\item \textbf{پهنای باند} (فقط در صورت انجام بخش ۳)\textbf{.} پهنای باند ۹۹ درصد توان را برای حالت‌های بخش
% [orig] \item \textbf{پهنای باند.} پهنای باند ۹۹ درصد توان را برای حالت‌های بخش
\ref{part:link}
و برای انحراف تیم حساب و با قاعدهٔ کارسون مقایسه کنید.
\item \textbf{کالیبراسیون انحراف} (فقط در صورت انجام بخش ۳)\textbf{.} نتیجهٔ روش صفرشدن حامل را با عدم‌قطعیت آن گزارش کنید.
% [orig] \item \textbf{کالیبراسیون انحراف.} نتیجهٔ روش صفرشدن حامل را با عدم‌قطعیت آن گزارش کنید.
% \item \textbf{اثر انحراف فرکانس.} نتایج چالش تیمی را تحلیل کنید: پهنای باند،
% \lr{SINAD}
% و آستانهٔ هر دو آشکارساز در برابر مقدار مرجع.
\end{enumerate}

% \subsection{تحلیل‌های الزامی جلسهٔ دوم}
% \begin{enumerate}
% \item \textbf{طیف باند پایهٔ استریو.} از فایل ضبط‌شدهٔ ایستگاه، سیگنال را با کد خودتان دمدوله و طیف باند پایه را رسم کنید؛ همهٔ مؤلفه‌ها را با فرکانس اندازه‌گیری‌شده علامت بزنید.
% \item \textbf{جداسازی استریو.} نتایج جدول
% \ref{tab:sep}
% را تحلیل کنید و اثر
% \code{lr_gain}
% و خطای تأخیر را با رابطهٔ خطای فاز بسنجید.
% \item \textbf{نرخ خطای دیکد \lr{RDS} در برابر سطح سیگنال.} منحنی نرخ خطای گروه را بر حسب
% $\mathrm{CNR}$
% رسم و
% $\mathrm{CNR}$
% متناظر با ۵۰ درصد را گزارش کنید؛ عدم‌قطعیت
% $\mathrm{CNR}$
% را برآورد کنید.
% \end{enumerate}

% \begin{question}[پرسش گسترشی]
% چرا در سیستم استریوی
% \lr{FM}
% از پایلوت
% \q{19}{kHz}
% استفاده می‌شود و نه خود حامل
% \q{38}{kHz}؟
% پاسخ را با استدلال طیفی و توان بدهید.
% \end{question}

\subsection{سایر اجزای الزامی}
\begin{itemize}
\item پاسخ پرسش‌های ۱ تا ۴ (پرسش‌های ۵ و ۶ اختیاری‌اند) همین دستور کار و جدول «پیش‌بینی در برابر اندازه‌گیری» هر جلسه با خطای نسبی و منبع کمّی هر خطا.
% [orig] \item پاسخ پرسش‌های ۱ تا ۶ (پرسش ۳ اختیاری است) همین دستور کار و جدول «پیش‌بینی در برابر اندازه‌گیری» هر جلسه با خطای نسبی و منبع کمّی هر خطا.
% [orig] \item پاسخ پرسش‌های ۱ تا ۶ همین دستور کار و جدول «پیش‌بینی در برابر اندازه‌گیری» هر جلسه با خطای نسبی و منبع کمّی هر خطا.
% [orig] \item پاسخ پرسش‌های ۱ تا ۱۹ همین دستور کار و جدول «پیش‌بینی در برابر اندازه‌گیری» هر جلسه با خطای نسبی و منبع کمّی هر خطا.
\item همان اجزایی که در آزمایش ۱ برشمرده شد: نمودار ساخته‌شده از دادهٔ خام با کد خودتان، بخش «آنچه اگر دوباره انجام می‌دادیم»، انتشار کد و فلوگراف‌ها و داده، خلاصهٔ اجرایی، جدول سهم اعضا و اعلام ابزارهای هوش مصنوعی. معیارهای ارزیابی همان جدول آیین‌نامه (بخش ۶) است.
\item پیشنهاد برای «عنصر تعاملی» وب‌سایت (آیین‌نامه، بخش ۵): پخش‌کنندهٔ صوتِ خروجی دمدولاسیون پیوند کابلی خودتان، پیش و پس از آستانه.
\end{itemize}

%======================================================================
\isection{bulb}{تمرین‌های تکمیلی (اختیاری)}
%======================================================================
\begin{enumerate}
% \item \textbf{بلوک‌های آماده.} خروجی گیرندهٔ خود را با بلوک
% \blk{WBFM Receive PLL}
% مقایسه کنید: جداسازی استریو و رفتار در سیگنال ضعیف.
\item \textbf{پیش‌تأکید.} در فرستنده
\blk{FM Preemphasis}
و در گیرنده
\blk{FM Deemphasis}
بگذارید و بهبود
$\mathrm{SNR}$
خروجی را برای تن
\q{10}{kHz}
اندازه بگیرید.
\item \textbf{دمدولاسیون آفلاین.} هر دو آشکارساز را در پایتون بنویسید و روی یکی از فایل‌های ضبط‌شده اجرا کنید؛ نتیجه را با خروجی
\GR{}
مقایسه کنید.
% \item \textbf{متن رادیویی.} پیام
% \lr{RadioText}
% یک ایستگاه یا فرستندهٔ کابلی را دیکد کنید و ساختار گروه حامل آن را شرح دهید.
\end{enumerate}

%======================================================================
\isection{list}{چک‌لیست پایان هر جلسه}
%======================================================================
\begin{checklist}
\item فلوگراف فرستنده متوقف شده، کابل‌ها جدا شده و هیچ فرستنده‌ای در حال کار نیست.
\item لاگ‌بوک جلسه با همهٔ مشاهدات، خطاها، مقدار
\code{f_fine}،
$P_{\mathrm{off}}$
و نقش اعضا در مخزن ثبت شده است.
\item همهٔ فایل‌های خام
\IQ{}
نام استاندارد دارند
(\code{lab02}،
نرخ
\val{240000})
و نسخهٔ پشتیبان آن‌ها گرفته شده است.
\item فلوگراف‌ها و تصاویر نمایشگر در مخزن است و جدول‌های همین جلسه کامل شده است.
\item هر دو
\lr{SDR}،
تضعیف‌گر و کابل‌ها (در صورت انجام بخش ۳) و آنتن شمارش و تحویل شده و برگهٔ تسویهٔ تجهیزات امضا شده است.
% [orig] تضعیف‌گر، بار ۵۰ اهمی، کابل‌ها و آنتن شمارش و تحویل شده و برگهٔ تسویهٔ تجهیزات امضا شده است.
\end{checklist}

%======================================================================
% \isection{flow}{پیوست: کد بلوک شمارندهٔ گروه \lr{RDS}}
%======================================================================
% این کد را در یک
% \blk{Python Block}
% قرار دهید (فایل
% \code{sdrlab_rds_counter.py}).
% بلوک فقط یک درگاه پیامی ورودی دارد و هر پیام بلوک
% \blk{RDS Decoder}
% را یک گروه درست می‌شمارد. کد بلوک ضبط همان کد پیوست آزمایش ۱ است.
% \begin{pycode}
% import threading
% import time

% import pmt
% from gnuradio import gr


% class blk(gr.basic_block):
%     def __init__(self, period=5.0):
%         gr.basic_block.__init__(self, name='SDR Lab RDS Group Counter',
%                                 in_sig=None, out_sig=None)
%         self.period = period
%         self._n = 0
%         self._run = False
%         self.message_port_register_in(pmt.intern('in'))
%         self.set_msg_handler(pmt.intern('in'), self._on_group)

%     def _on_group(self, msg):
%         self._n += 1          # one message = one correctly received group

%     def _loop(self):
%         while self._run:
%             time.sleep(self.period)
%             n, self._n = self._n, 0
%             nominal = 1187.5 / 104 * self.period
%             print('[RDS] %3d groups in %.0f s (nominal %.0f) -> group error rate %5.1f %%'
%                   % (n, self.period, nominal, 100 * max(0.0, 1 - n / nominal)))

%     def start(self):
%         self._run = True
%         threading.Thread(target=self._loop, daemon=True).start()
%         return True

%     def stop(self):
%         self._run = False
%         return True
% \end{pycode}

%======================================================================
\isection{book}{مراجع}
%======================================================================
\begin{latin}\small
\begin{enumerate}[label={[\arabic*]},leftmargin=2.2em,itemsep=1pt]
\item J. G. Proakis and M. Salehi, \textit{Communication Systems Engineering}, 2nd ed. Upper Saddle River, NJ: Prentice Hall, 2002.
\item M. Lichtman, \textit{PySDR: A Guide to SDR and DSP using Python}. [Online]. Available: https://pysdr.org
\item ITU-R, Recommendation BS.450, ``Transmission standards for FM sound broadcasting at VHF''; and Recommendation BS.643, ``Radio data system for automatic tuning and other applications in FM radio receivers for use with pilot-tone system.''
% \item B. Bloessl, \textit{gr-rds: FM RDS/TMC transceiver for GNU Radio}. [Online]. Available: https://github.com/bastibl/gr-rds
\item GNU Radio project, GNU Radio Wiki: block documentation. [Online]. Available: https://wiki.gnuradio.org
\end{enumerate}
\end{latin}

\end{document}
